\chapter{材料响应与辐射运动学}\label{chap:media-radiation}

前面求电子动力学时，曾把模式场和耦合矩阵元作为已知量。在具体实验中，这些量还需要由材料性质和装置几何求出。例如，同一个入射电子掠过金属表面或穿过介质时，会激发不同的场；即使总能量损失相同，能够传播到远处并被探测器收集的光也可能不同。本章从 Maxwell 方程求环境的响应，再把它代回电子的跃迁与探测公式。

先回忆线性方程的 Green 函数方法。求出点源产生的响应以后，任意分布源产生的场就等于各点响应的加权叠加。电流和电场都是矢量，源沿一个方向振荡时，产生的场可以沿另一个方向，因此这里需要一个带两个分量指标的 Green 张量。\begin{definition}[局域介质的推迟 Green 张量]
对各向同性、空间局域的非磁性介质，推迟 Green 张量 $\mathbf G^R$ 满足
\begin{equation}
{
\left[\nabla\times\nabla\times
-\frac{\omega^2}{c^2}\varepsilon(\vb{r},\omega)\right]
\mathbf G^R(\vb{r},\vb{r}';\omega)
=\mathbf I\,\delta^{(3)}(\vb{r}-\vb{r}').}
\label{eq:retarded_green_tensor_beginner_definition_dp123}
\end{equation}
其中 $\varepsilon$ 是相对介电函数，$\mathbf I$ 是三维单位张量。$\vb{r}'$ 是源点，$\vb{r}$ 是观察点，微分只作用于 $\vb{r}$。分量 $G^R_{ij}$ 记录沿 $j$ 方向的点源对 $i$ 方向电场的贡献。上标 $R$ 表示推迟边界条件：时域中的响应不能早于源，开放空间中的波向外传播。
\end{definition}

\begin{proposition}[源到场的线性映射]
取时间依赖为 $\ee^{-\ii\omega t}$ 的单频电流复振幅 $\vb{J}(\vb{r},\omega)$，相应电场复振幅为
\begin{equation}
{
\vb{E}(\vb{r},\omega)
=\ii\mu_0\omega\int\dd^3r'\,
\mathbf G^R(\vb{r},\vb{r}';\omega)\cdot
\vb{J}(\vb{r}',\omega).}
\label{eq:green_tensor_source_to_field_dp123}
\end{equation}
点号对源的方向指标求和，积分对所有源位置求和。这里 $[\mathbf G^R]={\rm m^{-1}}$；与 $\mu_0\omega$、电流密度和体积元相乘后，右边具有电场单位。对于一般时间信号，同一关系作用于逐频 Fourier 分量。改变材料或边界时，需要重求 $\mathbf G^R$；改变入射电子态时，需要重求作为源的电流。
\end{proposition}
\begin{proof}
令 $\mathcal L_\omega=\nabla\times\nabla\times-(\omega^2/c^2)\varepsilon$。
采用 $\ee^{-\ii\omega t}$ 约定，Maxwell 方程为
\[
 \nabla\times\vb{E}=\ii\omega\vb{B},
 \qquad
 \nabla\times\vb{B}=\mu_0\vb{J}-\ii\omega\mu_0\varepsilon_0\varepsilon\vb{E}.
\]
消去 $\vb{B}$ 得 $\mathcal L_\omega\vb{E}=\ii\mu_0\omega\vb{J}$。
对光滑紧支撑源，Green 张量的定义允许逐分量作用 $\mathcal L_\omega$：
\begin{align*}
 \mathcal L_\omega\left[\ii\mu_0\omega\int\dd^3r'\,
 \mathbf G^R(\vb{r},\vb{r}';\omega)\cdot\vb{J}(\vb{r}',\omega)\right]
 &=\ii\mu_0\omega\int\dd^3r'\,
 \mathbf I\delta^{(3)}(\vb{r}-\vb{r}')\cdot\vb{J}(\vb{r}',\omega)\\
 &=\ii\mu_0\omega\vb{J}(\vb{r},\omega).
\end{align*}
卷积继承 $\mathbf G^R$ 的推迟边界条件。在推迟预解式存在的复频率域，
具有同一源与同一边界条件的解唯一，故该卷积就是电场。
实频率上的辐射条件按 $\omega+\ii0^+$ 的边界值解释；本征极点处应以分布或极限吸收意义理解，
不能直接把奇异 Green 核当作有限矩阵。
更一般的源由该连续核在相应函数空间或分布空间中的延拓处理。
\end{proof}


电子源由量子态决定。若要分辨电子的初态 $|i\rangle$ 和末态 $|f\rangle$，使用电流算符在两态之间的矩阵元
\[
j^\mu_{fi}(x)=\langle f|\hat j^\mu(x)|i\rangle.
\]
这个跃迁电流保留电子反冲、自旋和波包相干。若电子波包足够窄、反冲不显著改变轨迹，并且所测量的量不分辨相关电子末态，才可以进一步使用沿给定轨迹运动的经典电流。\chapref{chap:electron-current-response}已经建立这一步近似；本章将在相应条件下用它求电子能量损失和辐射。\cite{GarciaDeAbajo2010RMP,GarciaDeAbajoDiGiulio2021,RoquesCarmes2023APR}

求出场以后，还要区分留在材料中的能量与进入远场的辐射，再考虑探测器实际接收的方向、频率和偏振。最后一节将把这些结果写成给定实验参数时的结果概率 $p(x|\vartheta)$，其中 $x$ 是探测结果，$\vartheta$ 表示待求的材料、近场或几何参数。由参数计算概率称为前向计算；利用测得的概率反过来估计参数，就是实验反演。

\section{辐射与介质响应}

电子在介质附近运动时，损失的能量可以转化为远场光，也可以留在材料的激发中。我们先从一般源电流求出远场辐射条件，再通过介电函数计算材料怎样响应。这样就能分别讨论加速、界面以及周期结构怎样使电子与可传播的光耦合，并在需要时保留电子反冲。

\subsection{加速与界面辐射}

匀速自由电子不能在真空中单独发射一个真实光子，因为初末电子和光子无法同时满足能量、动量守恒。外力可以向电子提供额外动量，界面或周期结构也可以参与交换，于是出现允许的辐射过程。为了统一计算这些情况，先写出任意守恒电流产生的远场，再分别代入相应的电子轨迹或诱导电流。

在进入点电荷轨迹以前，先保留任意守恒源。对正频率分量 $\vb{J}(\vb{r},\omega)$，真空远区场由推迟 Green 张量的辐射渐近式给出
\begin{equation}
{
\vb{E}(\vb{r},\omega)
=
\ii\mu_0\omega\frac{\ee^{\ii k r}}{4\pi r}
\bigl(\mathbf I-\vb{n}\vb{n}^{\mathsf T}\bigr)
\cdot\widetilde{\vb{J}}(k\vb{n},\omega)
+O(r^{-2}),
\qquad k=\omega/c.}
\label{eq:general_current_far_field_radiation}
\end{equation}
其中 $\vb{n}=\vb{r}/r$ 是观察方向，$\widetilde{\vb{J}}$ 是电流对空间坐标的 Fourier 变换。矩阵 $\mathbf I-\vb{n}\vb{n}^{\mathsf T}$ 去掉沿观察方向的分量，只留下横向电流；自变量 $k\vb{n}$ 则选出满足真空波色散 $k=\omega/c$ 的部分。因此，求某个过程的辐射，就是先求该过程的电流，再取它在传播波矢处的横向分量。对远场振幅取模平方并使用 Parseval 恒等式，即可求每单位频率和立体角的辐射能量。

点电荷轨迹提供这一一般源公式的第一个具体实现。对轨迹 $\vb{r}_e(t)$，四流由
\begin{equation*}
\rho(\vb{r},t)=q_{\mathrm e}\delta^{(3)}[\vb{r}-\vb{r}_e(t)],
\qquad
\vb{J}(\vb{r},t)=q_{\mathrm e}\vb{v}(t)\delta^{(3)}[\vb{r}-\vb{r}_e(t)]
\end{equation*}
定义。观察点 $\vb{r}$ 与推迟时间$t_{\rm r}$ 由
\begin{equation*}
t-t_{\rm r}=\frac{R(t_{\rm r})}{c},
\qquad
\vb{R}=\vb{r}-\vb{r}_e(t_{\rm r}),
\qquad
\vb{n}=\frac{\vb{R}}{R}
\end{equation*}
联系。相对论前向聚束可用下列无量纲量表示：
\begin{equation*}
\vb{\beta}_{\rm rel}=\frac{\vb{v}}{c},
\qquad
\gamma_{\rm L}=(1-\beta_{\rm rel}^2)^{-1/2},
\qquad
\mathfrak d_{\rm ret}=1-\vb{n}\cdot\vb{\beta}_{\rm rel}.
\end{equation*}
$\mathfrak d_{\rm ret}$ 是推迟时间的雅可比因子，也是相对论前向聚束的核心分母。由点电荷推迟势求得的远区 $R^{-1}$ 加速度场是点电荷因果电磁场的辐射部分；这一完整点电荷推迟场通常称为 Liénard--Wiechert 场。其远区项为
\begin{equation}
{
\vb{E}_{\rm rad}
=\frac{q_{\mathrm e}}{4\pi\varepsilon_0cR}
\frac{\vb{n}\times[(\vb{n}-\vb{\beta}_{\rm rel})\times\dot{\vb{\beta}}_{\rm rel}]}
{(1-\vb{n}\cdot\vb{\beta}_{\rm rel})^3}.}
\label{eq:lienard_wiechert_radiation_field}
\end{equation}
把这一场代入远区 Poynting 通量，并按源的推迟时间计量能流，得到角功率
\begin{equation}
{
\frac{\dd P_{\rm emit}}{\dd\Omega}
=
\frac{q_{\mathrm e}^2}{16\pi^2\varepsilon_0c}
\frac{
\left|
\vb{n}\times[(\vb{n}-\vb{\beta}_{\rm rel})
\times\dot{\vb{\beta}}_{\rm rel}]
\right|^2}
{(1-\vb{n}\cdot\vb{\beta}_{\rm rel})^5}.}
\label{eq:lienard_angular_emission_power}
\end{equation}
其中 $\dot{\vb{\beta}}_{\rm rel}=\vb{a}/c$。平行加速时角核正比于
\begin{equation*}
\frac{\sin^2\theta}{(1-\beta_{\rm rel}\cos\theta)^5},
\end{equation*}
因此 $\gamma_{\rm L}$ 增大时辐射逐渐压缩到前向角域。
\begin{figure}[htbp]
\centering
\includegraphics[width=0.74\linewidth]{lienard_angular_kernel.png}
\caption{Liénard 角功率在平行加速度下的归一化角核。曲线对应 $\beta_{\rm rel}=0.30,0.70,0.95$；相对论极限中的窄前向主瓣来自 $1-\vb{n}\cdot\vb{\beta}_{\rm rel}$的推迟时间的雅可比因子。}
\label{fig:lienard_angular_kernel}
\end{figure}
对全立体角积分后得到总辐射功率
\begin{equation}
{
P_{\rm rad}
=\frac{q_{\mathrm e}^2\gamma_{\rm L}^6}{6\pi\varepsilon_0c^3}
\left[a^2-|\vb{\beta}_{\rm rel}\times\vb{a}|^2\right].}
\label{eq:lienard_total_power_from_angular_flux}
\end{equation}
这给出了制动辐射的经典极限：外力首先改变电子四流，电子加速度的傅里叶分量随后耦合到真空传播模。量子论中“散射前发射”和“散射后发射”的振幅之和在软光子极限必须回到这一结果。

界面辐射的机制不同：电子可以保持近似匀速，但界面要求其 Coulomb/极化场在两侧重新满足 Maxwell 边界条件。为了分离这个几何机制，考虑速度在时间窗 $\Delta t$ 内由 $\vb{\beta}_i$ 变为 $\vb{\beta}_f$ 的有效端点，并定义
\begin{equation*}
\vb{T}(\vb{\beta}_{\rm rel})
\equiv
\frac{\vb{n}\times(\vb{n}\times\vb{\beta}_{\rm rel})}
{1-\vb{n}\cdot\vb{\beta}_{\rm rel}}.
\end{equation*}
当观测频率不能分辨跳变内部结构，即
\begin{equation*}
\omega(1-\vb{n}\cdot\vb{\beta}_{\rm rel})\Delta t\ll1,
\end{equation*}
每个有效端点只贡献 $\Delta\vb{T}$。多个相干端点在线性振幅层面相加，因此频谱角能量为
\begin{equation}
{
\frac{\dd^2 W}{\dd\omega\,\dd\Omega}
=\frac{1}{16\pi^3\varepsilon_0c}
\left|
\sum_jq_j
[\vb{T}(\vb{\beta}_{f,j})-\vb{T}(\vb{\beta}_{i,j})]
\ee^{\ii\varphi_j}
\right|^2.}
\label{eq:abrupt_velocity_change_spectrum}
\end{equation}
\begin{exbox}[理想导体渡越辐射]
设电荷 $q_{\mathrm e}$ 以速度 $v=\beta_{\rm rel}c$ 垂直入射理想导体平面。用真实电荷和镜像电荷满足真空侧边界条件，求每单位角频率和立体角的辐射能量。

\textbf{解。} 理想导体正入射可以用真实电荷和镜像电荷两个相干端点实现边界条件。此时得到理想导体正入射的经典跃迁辐射极限；这一结果通常称为 Ginzburg--Frank 公式：
\begin{equation}
{
\frac{\dd^2 W_{\rm TR}}{\dd\omega\,\dd\Omega}
=\frac{q_{\mathrm e}^2}{4\pi^3\varepsilon_0c}
\frac{\beta_{\rm rel}^2\sin^2\theta}
{(1-\beta_{\rm rel}^2\cos^2\theta)^2}.}
\label{eq:ginzburg_frank_perfect_conductor}
\end{equation}
这个理想化结果没有材料高频截止；真实高频抑制来自材料色散、有限界面厚度/几何和最终不可忽略的量子反冲，而不是单独由“形成长度”自动产生。

\medskip

理想导体正入射时，真空侧边界条件可由真实电荷与镜像电荷的两个相干端点实现。令真实电荷以 $+\beta_{\rm rel}\vb{e}_z$ 入射，镜像电荷 $-q_{\mathrm e}$ 以 $-\beta_{\rm rel}\vb{e}_z$ 镜像运动，并定义
\begin{equation*}
\vb{N}=\vb{n}\times(\vb{n}\times\vb{e}_z),
\qquad |\vb{N}|^2=\sin^2\theta.
\end{equation*}
代入\eqnrefs{eq:abrupt_velocity_change_spectrum} 的端点振幅核，两个贡献相加为
\begin{align*}
\mathcal A_{\rm TR}
&=-q_{\mathrm e}\frac{\beta_{\rm rel}\vb{N}}{1-\beta_{\rm rel}\cos\theta}
-q_{\mathrm e}\frac{\beta_{\rm rel}\vb{N}}{1+\beta_{\rm rel}\cos\theta}\\
&=-\frac{2q_{\mathrm e}\beta_{\rm rel}\vb{N}}{1-\beta_{\rm rel}^2\cos^2\theta}.
\end{align*}
取模平方并代回一般表达式，即得到\eqnrefs{eq:ginzburg_frank_perfect_conductor}。
\end{exbox}

上面的端点公式描述了局域辐射振幅，但有限长度结构还需要回答另一个问题：轨迹上相隔多远的两个发射位置仍会相干相加。为此定义\emph{形成长度}为纵向相位失配累积到给定量级所需的距离。取纵向失配相位累积到 $1\,\mathrm{rad}$ 的距离定义形成长度。它衡量相干积累距离，与材料中的真实衰减长度是不同的物理尺度。若某一频率分量沿轨迹的相位为
\begin{equation*}
\varphi_{\rm acc}(z)=\frac{\omega z}{v}-\kappa_zz,
\end{equation*}
则定义纵向失配 $\Delta \kappa_z=\omega/v-\kappa_z$，并以积累 $1\,\mathrm{rad}$ 相位为约定定义
\begin{equation}
{
L_f(\omega)\equiv\frac1{|\Delta \kappa_z|}
=\frac1{|\omega/v-\kappa_z|}.}
\label{eq:formation_length_general_mismatch}
\end{equation}
在真空远场 $\kappa_z=(\omega/c)\cos\theta$，因此精确形成长度为
\begin{equation*}
L_f
=\frac{c}{\omega\,[\beta_{\rm rel}^{-1}-\cos\theta]}
\qquad (0\le\theta<\pi,\ \beta_{\rm rel}<1).
\end{equation*}
为控制超相对论小角展开，定义
\begin{equation*}
 x\equiv\gamma_{\rm L}^{-2},
 \qquad y\equiv\theta^2,
\end{equation*}
以及两个精确余项
\begin{align*}
 r_\beta(x)
 &\equiv(1-x)^{-1/2}-1-\frac{x}{2},\\
 r_\theta(\theta)
 &\equiv1-\cos\theta-\frac{\theta^2}{2}.
\end{align*}
于是
\begin{equation*}
\beta_{\rm rel}^{-1}-\cos\theta
=\frac{x+y}{2}+r_\beta+r_\theta.
\end{equation*}
把
\begin{equation*}
\delta_f\equiv\frac{2(r_\beta+r_\theta)}{x+y}
\end{equation*}
定义为分母的相对修正，便得到精确关系
\begin{equation}
{
L_f
=\frac{2c}{\omega(x+y)}\frac{1}{1+\delta_f},
\qquad
x=\gamma_{\rm L}^{-2},\ y=\theta^2.}
\label{eq:formation_length_ultrarelativistic}
\end{equation}
节末推导给出显式上界 $|\delta_f|\le\epsilon_f$。当 $\epsilon_f<1$ 时，领先小角式
$L_f^{(0)}=2c/[\omega(\gamma_{\rm L}^{-2}+\theta^2)]$
满足
\begin{equation*}
\left|\frac{L_f-L_f^{(0)}}{L_f^{(0)}}\right|
\le\frac{\epsilon_f}{1-\epsilon_f}.
\end{equation*}
因此是否可以使用常见的小角形成长度，不再由“Lorentz 因子很大、$\theta$ 很小”的定性判断决定，而由这个显式误差界决定。
周期性或可分辨反冲会改变纵向相位匹配，此时精确失配写成
\begin{equation*}
\Delta \kappa_z
=\frac{p_{i,z}-p_{f,z}}{\hbar}-\kappa_z-G_z.
\end{equation*}
电子真实初末动量、辐射波矢和倒格矢因此共同进入形成长度；界面辐射、周期辐射与有限长度同步都可据此归结为纵向相位在何种距离上失去相干。

对有损或开放纳米结构，还必须区分“电子失去的总能量”和“真正跑到远场探测器的辐射能量”。EELS 读取前者；阴极发光（CL）或远场光子计数读取后者中被光学系统实际收集的部分。二者的通道分解直接来自同一个频域 Poynting 定理。采用全文的 $\ee^{-\ii\omega t}$ 约定，对给定经典频域电流 $\widetilde{\vb{J}}$，Green 张量给
\[
\widetilde{\vb{E}}(\vb{r},\omega)
=\ii\mu_0\omega\int\dd^3r'\,
\mathbf G^R(\vb{r},\vb{r}';\omega)\cdot
\widetilde{\vb{J}}(\vb{r}',\omega).
\]
源在该频率上损失给电磁环境的平均功率定义为
\begin{equation}
{
P_{\rm loss}(\omega)
\equiv
-\frac12\Re\int\dd^3r\,
\widetilde{\vb{J}}^{*}\cdot\widetilde{\vb{E}}
=\frac{\mu_0\omega}{2}
\iint\dd^3r\,\dd^3r'\,
\widetilde{\vb{J}}^{*}(\vb{r})\cdot
\operatorname{Im}\mathbf G^R(\vb{r},\vb{r}';\omega)
\cdot\widetilde{\vb{J}}(\vb{r}').}
\label{eq:source_total_loss_green_power}
\end{equation}
对于空间局域、非磁被动介质，这个总功率严格分成材料吸收与无穷远辐射通量：
\begin{equation}
{
P_{\rm loss}=P_{\rm abs}+P_{\rm rad},
\qquad
P_{\rm abs}
=\frac{\omega\varepsilon_0}{2}
\int_V\dd^3r\,
\widetilde{\vb{E}}^{*}\cdot
\boldsymbol\varepsilon_I\cdot
\widetilde{\vb{E}},
\qquad
P_{\rm rad}
=\frac12\Re\oint_{\partial V}\dd\vb{A}\cdot
(\widetilde{\vb{E}}\times\widetilde{\vb{H}}^{*}).}
\label{eq:loss_absorption_radiation_decomposition}
\end{equation}
$V$ 取到包围全部材料并延伸到辐射远区。若还存在磁损耗，应在 $P_{\rm abs}$ 中增加
$(\omega\mu_0/2)\int \vb{H}^*\cdot\boldsymbol\mu_I\cdot\vb{H}$。因此 $\operatorname{Im}G^R$ 本身包含所有开放通道：材料耗散和远场辐射只有在作 Poynting 分解后才被区分。

这给出自由电子实验中一个非常实用的层级关系。经典 EELS 测得的是电子的总能量损失，原则上对应 $P_{\rm loss}$；阴极发光或远场光子计数只读取 $P_{\rm rad}$ 中进一步被数值孔径、偏振和频带选中的部分。因此可以定义
\begin{equation}
{
\mathcal B_{\rm rad}(\omega)
\equiv\frac{P_{\rm rad}(\omega)}{P_{\rm loss}(\omega)},
\qquad 0\le\mathcal B_{\rm rad}\le1,}
\label{eq:radiative_branching_efficiency_green}
\end{equation}
但 $\mathcal B_{\rm rad}$ 是环境通道分支比，不是电子--环境微观耦合常数。量子跃迁流 $\vb{j}_{fi}$ 的一量子过程也具有完全相同的通道分解：先由 $\operatorname{Im}G^R$ 给总环境谱权重，再由远场投影或材料吸收算符选择具体观测通道。

\subsection{材料响应参数化}

Green 张量回答“已知材料响应后，给定源产生什么场”，但材料为什么极化仍需由本构动力学给出。在线性范围内，外电场驱动带电自由度；惯性、回复力、阻尼以及空间压力决定它们的位移，位移再组成宏观极化 $\vb{P}$。解掉这些线性运动方程便得到 $\vb{P}$ 对 $\vb{E}$ 的响应核，等价地得到介电函数 $\boldsymbol\varepsilon_r(\vb{\kappa},\omega)$。

因此 Lorentz、Drude 与流体模型并不是三条互不相关的经验公式：Lorentz 振子保留回复力，Drude 模型取回复力为零，流体模型再加入空间梯度压力而产生 $\vb{\kappa}$ 依赖。材料模型改变的是 Maxwell 算符中的本构核；电子跃迁流与 Green 张量耦合的结构不变。介电函数的极点和零点控制材料共振，虚部编码耗散，显式波矢依赖则量化空间色散。


先考虑空间局域的束缚电荷 Lorentz 振子。令
\begin{equation*}
n_{\rm osc}\ [{\rm m^{-3}}],
\quad
m_{\rm osc}\ [{\rm kg}],
\quad
q_{\rm osc}\ [{\rm C}],
\quad
\omega_0,\Gamma_{\rm mat}\ [{\rm s^{-1}}]
\end{equation*}
分别表示振子数密度、单振子质量、电荷、本征角频率和阻尼率；位移 $x(t)$ 的单位为米。线性运动方程为
\begin{equation*}
m_{\rm osc}\ddot x
+m_{\rm osc}\Gamma_{\rm mat}\dot x
+m_{\rm osc}\omega_0^2x
=q_{\rm osc}E.
\end{equation*}
定义振子强度频率
\begin{equation*}
\Omega_{\rm osc}^2
\equiv
\frac{n_{\rm osc}q_{\rm osc}^2}
{\varepsilon_0m_{\rm osc}},
\qquad
[\Omega_{\rm osc}]={\rm s^{-1}},
\end{equation*}
则频域极化率为
\begin{equation*}
\chi_{\rm L}(\omega)
=\frac{\Omega_{\rm osc}^2}
{\omega_0^2-\omega^2-\ii\Gamma_{\rm mat}\omega}.
\end{equation*}
分子与分母都具有 ${\rm s^{-2}}$，所以 $\chi_{\rm L}$ 无量纲。

自由载流子对应 $\omega_0=0$。对电子气体令平衡数密度为 $n_0$，定义等离子体角频率
\begin{equation*}
\omega_p^2
\equiv\frac{n_0e^2}{\varepsilon_0m_{\mathrm e}},
\qquad
[\omega_p]={\rm s^{-1}}.
\end{equation*}
若高频束缚电子背景用无量纲常数 $\varepsilon_\infty$ 表示，则局域 Drude 介电函数为
\begin{equation*}
\varepsilon_D(\omega)
=\varepsilon_\infty
-\frac{\omega_p^2}
{\omega(\omega+\ii\Gamma_{\rm mat})}.
\end{equation*}
$\Gamma_{\rm mat}^{-1}$ 是动量弛豫时间尺度；在被动约定 $\ee^{-\ii\omega t}$ 下，$\Gamma_{\rm mat}>0$ 把极点推入复频率下半平面。

局域 Drude 模型忽略电子流体的压缩能，因此材料响应只依赖 $\omega$。当空间变化尺度接近电子气体内部压缩长度时，纵向压力梯度必须进入电子流体的运动方程。令 $\delta n(\vb{r},t)$ 为相对 $n_0$ 的数密度涨落，单位 ${\rm m^{-3}}$；$\vb{u}(\vb{r},t)$ 为流体速度，单位 ${\rm m\,s^{-1}}$；$v_H$ 为表征压缩刚度的速度尺度，单位 ${\rm m\,s^{-1}}$。线性化连续性方程和动量方程为
\begin{equation*}
\partial_t\delta n+n_0\nabla\cdot\vb{u}=0,
\end{equation*}
\begin{equation*}
(\partial_t+\Gamma_{\rm mat})\vb{u}
=\frac{q_{\mathrm e}}{m_{\mathrm e}}\vb{E}
-\frac{v_H^2}{n_0}\nabla\delta n.
\end{equation*}
对纵向模取
\[
\delta n(\vb{r},t)=\delta n_0\ee^{\ii\vb{\kappa}\cdot\vb{r}-\ii\omega t},\qquad
\vb{u}(\vb{r},t)=u_0\hat{\vb{\kappa}}\ee^{\ii\vb{\kappa}\cdot\vb{r}-\ii\omega t},
\]
其中 $\kappa=|\vb{\kappa}|$ 的单位为 ${\rm m^{-1}}$。消去 $\delta n_0$ 和 $u_0$ 后得到纵向非局域介电函数
\begin{equation*}
\varepsilon_L(\kappa,\omega)
=\varepsilon_\infty
-\frac{\omega_p^2}
{\omega(\omega+\ii\Gamma_{\rm mat})-v_H^2\kappa^2}.
\end{equation*}
这里空间色散由无量纲比值
\begin{equation}
{
\epsilon_{\rm nl}(\kappa,\omega)
\equiv
\frac{v_H^2\kappa^2}
{|\omega(\omega+\ii\Gamma_{\rm mat})|}}
\label{eq:hydrodynamic_nonlocal_control}
\end{equation}
控制。$\epsilon_{\rm nl}\ll1$ 时，压力梯度相对局域惯性项较小；当 $\epsilon_{\rm nl}$ 达到一阶量级时，$\kappa$ 依赖不能再从本构关系中删去。

无外部自由电荷的纵向体模满足 $D_L=0$，即 $\varepsilon_L(\kappa,\omega)=0$。在 $\Gamma_{\rm mat}=0$ 时，正频支为
\begin{equation}
{
\omega_L^2(\kappa)
=\frac{\omega_p^2}{\varepsilon_\infty}
+v_H^2\kappa^2.}
\label{eq:hydrodynamic_bulk_longitudinal_dispersion}
\end{equation}
第一项是局域体等离子体频率，第二项是压缩性造成的空间色散。

色散介质的能量密度也由同一显式材料自由度决定。对无损 Lorentz 振子，周期平均机械能与电场耦合能消去振子坐标后，可把电场与物质部分写成
\begin{equation*}
\bar u_{E+{\rm mat}}
=\frac{\varepsilon_0}{4}
\frac{\partial[\omega\varepsilon(\omega)]}{\partial\omega}
|E|^2.
\end{equation*}
$\partial_\omega[\omega\varepsilon]$ 无量纲，因此该式单位为 ${\rm J\,m^{-3}}$。完整周期平均能量还包含磁场项 $\mu_0|H|^2/4$。这个频率导数记录的是被消去材料振子的储能贡献，因此在强色散区必须与 Maxwell 场能一起归一化。

\medskip

\eqnrefs{eq:hydrodynamic_nonlocal_control} 已把压强梯度引入电子流体，从而使纵向响应依赖波数。将连续性方程和纵向动量方程联立，消去密度涨落与流速，并要求无外部自由电荷时存在非零纵向电场，就得到\eqnrefs{eq:hydrodynamic_bulk_longitudinal_dispersion}。对流体动力学模型取纵向平面波。连续性方程变成
\[
-\ii\omega\delta n
+\ii n_0\kappa u_L=0,
\]
因此
\[
\delta n=\frac{n_0\kappa}{\omega}u_L.
\]
纵向动量方程为
\[
(-\ii\omega+\Gamma_{\rm mat})u_L
=\frac{q_{\mathrm e}}{m_{\mathrm e}}E_L
-\ii\frac{v_H^2\kappa}{n_0}\delta n.
\]
代入 $\delta n$：
\[
(-\ii\omega+\Gamma_{\rm mat})u_L
=\frac{q_{\mathrm e}}{m_{\mathrm e}}E_L
-\ii\frac{v_H^2\kappa^2}{\omega}u_L.
\]
两边乘 $\ii\omega$：
\[
[\omega(\omega+\ii\Gamma_{\rm mat})-v_H^2\kappa^2]u_L
=\ii\omega\frac{q_{\mathrm e}}{m_{\mathrm e}}E_L.
\]
因此
\[
u_L
=\frac{\ii\omega q_{\mathrm e}/m_{\mathrm e}}
{\omega(\omega+\ii\Gamma_{\rm mat})-v_H^2\kappa^2}E_L.
\]
自由电子电流为 $J_L=q_{\mathrm e} n_0u_L$，而极化电流满足 $J_L=-\ii\omega P_L$，所以
\begin{align*}
P_L
&=\frac{\ii}{\omega}q_{\mathrm e} n_0u_L\\
&=-\varepsilon_0
\frac{\omega_p^2}
{\omega(\omega+\ii\Gamma_{\rm mat})-v_H^2\kappa^2}E_L.
\end{align*}
总纵向位移场
\[
D_L=\varepsilon_0\varepsilon_\infty E_L+P_L
\]
因此
\[
\varepsilon_L(\kappa,\omega)
=\varepsilon_\infty-
\frac{\omega_p^2}
{\omega(\omega+\ii\Gamma_{\rm mat})-v_H^2\kappa^2}.
\]
无外部自由电荷时 $D_L=0$；非零 $E_L$ 要求 $\varepsilon_L=0$。取 $\Gamma_{\rm mat}=0$：
\begin{align*}
0
&=\varepsilon_\infty-
\frac{\omega_p^2}{\omega^2-v_H^2\kappa^2},\\
\varepsilon_\infty(\omega^2-v_H^2\kappa^2)
&=\omega_p^2,\\
\omega^2
&=\frac{\omega_p^2}{\varepsilon_\infty}+v_H^2\kappa^2,
\end{align*}
得到\eqnrefs{eq:hydrodynamic_bulk_longitudinal_dispersion}。


均匀各向同性但具有空间色散的介质还可以做一个非常重要的结构分解。对给定非零波矢 $\vb{\kappa}$，定义纵向与横向投影算符
\begin{equation}
{
\mathsf P_L(\vb{\kappa})
=\hat{\vb{\kappa}}\hat{\vb{\kappa}}^{\mathsf T},
\qquad
\mathsf P_T(\vb{\kappa})
=\mathsf I-\mathsf P_L(\vb{\kappa}).}
\label{eq:nonlocal_LT_projectors}
\end{equation}
它们满足
$\mathsf P_{L,T}^2=\mathsf P_{L,T}$、
$\mathsf P_L\mathsf P_T=0$、
$\mathsf P_L+\mathsf P_T=I$。旋转不变性要求任意二阶响应张量只能在这两个不变子空间上取标量值，因此最一般的均匀各向同性非局域介电张量不是一个标量 $\varepsilon(\kappa,\omega)$，而是
\begin{equation}
{
\boldsymbol\varepsilon(\vb{\kappa},\omega)
=\varepsilon_T(\kappa,\omega)\mathsf P_T
+\varepsilon_L(\kappa,\omega)\mathsf P_L.}
\label{eq:isotropic_nonlocal_dielectric_LT_decomposition}
\end{equation}
$\varepsilon_T$ 控制横向电磁/极化激元分支，$\varepsilon_L$ 控制电荷密度和纵向等离激元分支；空间局域各向同性介质才进一步满足
$\varepsilon_T=\varepsilon_L=\varepsilon(\omega)$。

把这个分解代入频域 Maxwell 算符。由于
\[
\ii\vb{\kappa}\times
(\ii\vb{\kappa}\times\vb{E})
=\kappa^2\mathsf P_T\vb{E},
\]
有
\begin{equation*}
\mathsf L(\vb{\kappa},\omega)
=\left[\kappa^2-\frac{\omega^2}{c^2}\varepsilon_T\right]\mathsf P_T
-\frac{\omega^2}{c^2}\varepsilon_L\mathsf P_L.
\end{equation*}
两个投影子空间彼此正交，所以逆算符可逐块求逆：
\begin{equation}
{
\mathbf G^R(\vb{\kappa},\omega)
=\frac{\mathsf P_T}
{\kappa^2-(\omega^2/c^2)\varepsilon_T(\kappa,\omega)}
-\frac{c^2}{\omega^2}
\frac{\mathsf P_L}{\varepsilon_L(\kappa,\omega)}.}
\label{eq:isotropic_nonlocal_green_LT_decomposition}
\end{equation}
因此环境谱的两类极点来源不同：横向传播模满足
\begin{equation*}
\kappa^2=\frac{\omega^2}{c^2}\varepsilon_T(\kappa,\omega),
\end{equation*}
而无外部自由电荷的纵向体模满足
\begin{equation*}
\varepsilon_L(\kappa,\omega)=0.
\end{equation*}
这正是前面流体模型的体等离激元色散出现于 $\varepsilon_L=0$、而普通光学折射主要读取 $\varepsilon_T$ 的原因。

自由电子的电流同样可分成
\[
\vb{J}=\vb{J}_T+\vb{J}_L,
\qquad
\vb{J}_{T,L}=\mathsf P_{T,L}\vb{J}.
\]
由于两个投影正交，均匀介质中的 Green 能损核严格分裂为
\begin{equation}
{
\begin{aligned}
\vb{J}^*\!\cdot\operatorname{Im}\mathbf G^R\!\cdot\vb{J}
={}&|\vb{J}_T|_{\operatorname{Im}G_T}^2
+|\vb{J}_L|_{\operatorname{Im}G_L}^2,
\end{aligned}}
\label{eq:electron_loss_LT_split}
\end{equation}
其中
\[
G_T=\frac1{\kappa^2-(\omega^2/c^2)\varepsilon_T},
\qquad
G_L=-\frac{c^2}{\omega^2\varepsilon_L},
\]
并约定
\[
|\vb{J}_T|_{\operatorname{Im}G_T}^2
\equiv
\vb{J}_T^*\!\cdot\operatorname{Im}(G_T\mathsf P_T)\cdot\vb{J}_T,
\qquad
|\vb{J}_L|_{\operatorname{Im}G_L}^2
\equiv
\vb{J}_L^*\!\cdot\operatorname{Im}(G_L\mathsf P_L)\cdot\vb{J}_L.
\]
第一项对应电子耦合到横向辐射/极化激元通道，第二项对应纵向电荷密度通道。真空中 $\varepsilon_L=1$ 且 $G_L$ 没有耗散极点，所以匀速电子不能把能量不可逆地送入“纵向真空光子”；材料中的 $\varepsilon_L=0$ 极点则允许真正的体等离激元损失。

\begin{proposition}[纵横投影的 Green 逆与正交二次型]
固定非零实波矢，令 $P_L=\vb{\kappa}\vb{\kappa}^{\mathsf T}/\kappa^2$、$P_T=I-P_L$。对无手性、互易、各向同性均匀响应 $\varepsilon=\varepsilon_LP_L+\varepsilon_TP_T$，若两分母非零，Maxwell 算符逆为\eqnrefs{eq:isotropic_nonlocal_green_LT_decomposition}；任意电流的 Green 二次型无纵横交叉项，给\eqnrefs{eq:electron_loss_LT_split}。极点按推迟边界值解释；仅有旋转不变性而允许手性时，还可能存在反对称横向响应。
\end{proposition}
\begin{proof}
固定 $\vb{\kappa}\neq0$。旋转不变性把响应分成纵向一维与横向二维两个互不混合的子空间，故
\[
\varepsilon
=\varepsilon_L\mathsf P_L+\varepsilon_T\mathsf P_T,
\]
得到正文\eqnrefs{eq:isotropic_nonlocal_dielectric_LT_decomposition}。

频域 Maxwell 方程采用全文 $\ee^{-\ii\omega t}$ 约定：
\[
\left[
\nabla\times\nabla\times
-\frac{\omega^2}{c^2}\boldsymbol\varepsilon
\right]\vb{E}
=\ii\mu_0\omega\vb{J}.
\]
对平面波 Fourier 分量 $\nabla\to\ii\vb{\kappa}$，有
\begin{align*}
\nabla\times\nabla\times\vb{E}
&=-\vb{\kappa}\times
(\vb{\kappa}\times\vb{E})\\
&=\kappa^2\vb{E}
-\vb{\kappa}(\vb{\kappa}\cdot\vb{E})\\
&=\kappa^2\mathsf P_T\vb{E}.
\end{align*}
因此
\[
\mathsf L
=\left(\kappa^2-k_0^2\varepsilon_T\right)\mathsf P_T
-k_0^2\varepsilon_L\mathsf P_L,
\qquad k_0=\omega/c.
\]
设逆核为 $G_T\mathsf P_T+G_L\mathsf P_L$。利用投影代数
$\mathsf P_T\mathsf P_L=0$，要求 $\mathsf L\mathbf G=I$ 给
\[
(\kappa^2-k_0^2\varepsilon_T)G_T=1,
\qquad
(-k_0^2\varepsilon_L)G_L=1,
\]
于是
\[
G_T=\frac1{\kappa^2-k_0^2\varepsilon_T},
\qquad
G_L=-\frac1{k_0^2\varepsilon_L},
\]
即正文\eqnrefs{eq:isotropic_nonlocal_green_LT_decomposition}。

最后把任意电流写成
$\vb{J}=\vb{J}_T+\vb{J}_L$。由于
$\vb{J}_T^*\mathsf P_L\vb{J}=0$、
$\vb{J}_L^*\mathsf P_T\vb{J}=0$，有
\begin{align*}
\vb{J}^*\cdot\operatorname{Im}\mathbf G\cdot\vb{J}
&=\vb{J}_T^*\cdot
\operatorname{Im}(G_T\mathsf P_T)\cdot\vb{J}_T\\
&\quad+\vb{J}_L^*\cdot
\operatorname{Im}(G_L\mathsf P_L)\cdot\vb{J}_L,
\end{align*}
没有纵横交叉项，从而得到\eqnrefs{eq:electron_loss_LT_split}。这个正交分解只依赖均匀各向同性结构；在界面处波矢法向分量不再是全局好量子数，必须重新求边值 Green 张量。
\end{proof}


纵横分解也说明了为什么强非局域体系不能只用一个折射率 $n(\omega)$ 描述。折射率编码的是横向传播质量壳；自由电子近场却可以携带远大于 $\omega/c$ 的波矢，因此 EELS 能直接采样大 $\kappa$ 的纵向响应。二维电子气、金属纳米结构和电子液体中的空间色散往往先在这种近场通道中显现。

式~\eqnrefs{eq:isotropic_nonlocal_green_LT_decomposition} 依赖“均匀 + 各向同性”这两个条件。界面、非均匀或各向异性结构不再由全局的 $\mathsf P_L/\mathsf P_T$ 同时对角化，此时应退回双点 Green 算符或相应表面响应核，再由具体对称性做进一步分解。

三维介电核仍不是低维材料最自然的描述。对原子级薄膜、二维电子气或表面导电层，若层厚远小于电子所采样的法向场变化尺度，把它硬塞进一个厚度依赖的体介电常数会引入人为参数；更普适的对象是定义在二维界面上的表面电导核。设界面位于 $z=0$，面内坐标记为 $\vb{\rho}$，则最一般线性表面响应写成
\begin{equation}
{
K_a(\vb{\rho},\omega)
=
\sum_{b=x,y}
\int\dd^2\rho'\,
\sigma_{ab}(\vb{\rho},\vb{\rho}';\omega)
E_{\parallel,b}(\vb{\rho}',0,\omega),}
\label{eq:general_2d_conductivity_kernel}
\end{equation}
其中 $\vb{K}$ 是表面电流密度，单位 ${\rm A\,m^{-1}}$；$\sigma_{ab}$ 是二维非局域电导核。只有面内平移不变时才可以 Fourier 对角化为
\[
\vb{K}(\vb{q},\omega)
=
\boldsymbol\sigma(\vb{q},\omega)
\vb{E}_\parallel(\vb{q},0,\omega).
\]
因此常写的标量面电导率 $\sigma(\omega)$ 已经同时使用了“局域 + 各向同性 + 均匀”的三重特例。

零厚度表面层并不改变 Maxwell 方程在两侧体介质中的形式，而是改变边界条件。取法向单位矢量 $\hat{\vb{n}}$ 从介质 $1$ 指向介质 $2$，则
\begin{equation}
{
\hat{\vb{n}}\times(\vb{E}_2-\vb{E}_1)=0,
\qquad
\hat{\vb{n}}\times(\vb{H}_2-\vb{H}_1)=\vb{K}.}
\label{eq:conductive_sheet_maxwell_jump}
\end{equation}
第一式保证切向电场连续，第二式说明表面电流就是切向磁场的跳跃。若周围两侧都是均匀各向同性体介质，固定面内波矢 $\vb{q}$ 后可用 TE/TM 基底。定义法向波数
\[
k_{zj}
=
\sqrt{\varepsilon_j(\omega)\frac{\omega^2}{c^2}-q^2},
\qquad
\operatorname{Im}k_{zj}\ge0,
\]
以及切向波导纳
\begin{equation*}
Y_{p,j}
\equiv
\frac{\varepsilon_0\varepsilon_j\omega}{k_{zj}},
\qquad
Y_{s,j}
\equiv
\frac{k_{zj}}{\mu_0\omega}.
\end{equation*}
对面内各向同性但允许空间色散的二维层，电导同样分解成
\begin{equation}
{
\boldsymbol\sigma(\vb{q},\omega)
=
\sigma_L(q,\omega)\mathsf P_L^\parallel
+
\sigma_T(q,\omega)\mathsf P_T^\parallel,}
\label{eq:2d_conductivity_LT_decomposition}
\end{equation}
其中
$\mathsf P_L^\parallel=\hat{\vb{q}}\hat{\vb{q}}^{\mathsf T}$，
$\mathsf P_T^\parallel=I_\parallel-\mathsf P_L^\parallel$。TM 场的切向电场平行于 $\vb{q}$，因此只读取 $\sigma_L$；TE 场的切向电场垂直于 $\vb{q}$，因此只读取 $\sigma_T$。

无外部入射场的表面本征模要求边界线性方程存在非零解。于是两个极化的极点条件分别为
\begin{equation}
{
D_p(q,\omega)
\equiv
Y_{p,1}+Y_{p,2}+\sigma_L(q,\omega)=0,
\qquad
D_s(q,\omega)
\equiv
Y_{s,1}+Y_{s,2}+\sigma_T(q,\omega)=0.}
\label{eq:2d_sheet_surface_mode_conditions}
\end{equation}
符号取决于全文的法向与 Fourier 约定，但零点条件本身与同一套边界约定自洽。对真正各向异性或 Hall 型二维层，TE/TM 不再独立；把两种切向偏振组成二维列向量以后，本征条件提升为
\begin{equation}
{
\det\!\left[
\mathbf Y_1(\vb{q},\omega)
+\mathbf Y_2(\vb{q},\omega)
+\boldsymbol\sigma(\vb{q},\omega)
\right]=0.}
\label{eq:2d_sheet_anisotropic_pole_determinant}
\end{equation}
这比给每一种二维材料单独写色散关系更一般：材料全部进入 $\boldsymbol\sigma$，几何与外部介质进入两侧导纳。

对自由电子而言，二维层的 EELS 不需要另造一套跃迁理论：仍从一般 Green 损失核出发，只把材料响应换成二维边界条件。平面平移对称性允许对 $\vb{\rho}-\vb{\rho}'$ 作 Fourier 变换，于是每个面内动量 $\vb{q}$ 独立散射；反射 Green 张量的极点与连续谱分别由\eqnrefs{eq:2d_sheet_surface_mode_conditions,eq:2d_sheet_anisotropic_pole_determinant} 控制。电子倏逝近场可以提供 $q\gg\omega/c$ 的大面内动量，而远场光子受光锥限制，因此 EELS 对 $\sigma_L(q,\omega)$ 的空间色散和纵向集体模尤其敏感\cite{Simonaitis2025EELS2D}。实验谱峰的位置主要读取边界预解式的极点，峰强和线形还包含电子轨迹投影、极点留数与耗散连续谱；这些因素在下一步显式 Green 张量中同时出现。

为求界面反射产生的电场，接下来把反射系数代入 Green 张量。对源点和观察点都位于介质 $1$ 的情形，固定 $(\vb{q},\omega)$ 后的散射 Green 张量具有 Weyl 形式
\begin{equation}
{
\mathbf G_{\rm sc}^R(\vb{q};z,z';\omega)
=\frac{\ii}{2k_{z1}}\ee^{\ii k_{z1}(z+z')}
\left[r_s\,\vb{e}_s^+\vb{e}_s^-+r_p\,\vb{e}_p^+\vb{e}_p^-\right].}
\label{eq:2d_sheet_reflected_green_weyl}
\end{equation}
采用标准 Fresnel 的 $p$ 极化振幅约定，各向同性非局域二维层的反射系数为
\begin{align}
{r_s
=\frac{Y_{s,1}-Y_{s,2}-\sigma_T}
{Y_{s,1}+Y_{s,2}+\sigma_T},}
\label{eq:2d_sheet_rs_conductivity}\\
{r_p
=\frac{Y_{p,2}-Y_{p,1}+\sigma_L}
{Y_{p,1}+Y_{p,2}+\sigma_L}.}
\label{eq:2d_sheet_rp_conductivity}
\end{align}
两个分母正好就是\eqnrefs{eq:2d_sheet_surface_mode_conditions} 的 $D_s,D_p$，所以 EELS 谱峰、表面本征模和 Green 张量极点是同一个边界预解式的不同读取方式。若改用切向电场本身而不是标准 $p$ 极化基矢定义反射振幅，$r_p$ 的整体符号会随基底约定变化；物理 Green 张量在同一约定下不变。

对平行界面的电子轨迹 $\vb{r}_e(t)=(vt,0,b)$，时间 Fourier 变换自动给
\[
q_x=\frac{\omega}{v}.
\]
把\eqnrefs{eq:2d_sheet_reflected_green_weyl} 代入一般 Green 损失核后，只剩横向面内动量 $q_y$ 的积分。倏逝区 $q>\sqrt{\varepsilon_1}\omega/c$ 中写 $k_{z1}=\ii\kappa_1$，则反射项在同一高度 $z=z'=b$ 取样时都带
\begin{equation}
{\ee^{-2\kappa_1b}.}
\label{eq:2d_eels_evanscent_impact_parameter_factor}
\end{equation}
所以冲击参数 $b$ 直接决定电子能够访问多大的面内波矢。小 $b$ 不只是“耦合更强”，而是扩大了对高 $q$ 非局域纵向响应 $\sigma_L(q,\omega)$ 的采样窗口。

平面界面的反射系数利用了平移对称性。对于有限颗粒或任意形状的介电结构，
可以先选取一个已知背景，再把真实结构与背景的差别写成散射算符。
这样，材料对比进入局部作用，背景 Green 核负责把每次产生的场传播出去。

\begin{proposition}[背景响应与 Maxwell 散射算符]
固定一个复频率，取场空间 $\mathcal F$、源空间 $\mathcal S$，
以及包含相同推迟边界条件的背景 Maxwell 实现
$\mathsf L_0:D(\mathsf L_0)\subset\mathcal F\to\mathcal S$。
设它的逆 $\mathsf G_0:\mathcal S\to\mathcal F$ 有界，
对比算符 $\mathsf V:\mathcal F\to\mathcal S$ 有界，
且 $I-\mathsf G_0\mathsf V$ 在 $\mathcal F$ 上有有界逆。
真实算符取 $\mathsf L=\mathsf L_0-\mathsf V$，定义域保持 $D(\mathsf L_0)$。
给定源 $f\in\mathcal S$，令背景场 $\vb{E}_{\rm inc}=\mathsf G_0f$，则真实场满足
\begin{equation}
\vb{E}=\vb{E}_{\rm inc}+\mathsf G_0\mathsf V\vb{E}.
\label{eq:feqo-maxwell-lippmann}
\end{equation}
Maxwell 散射算符及真实逆为
\begin{equation}
\mathsf T=\mathsf V(I-\mathsf G_0\mathsf V)^{-1},\qquad
\mathsf G=\mathsf G_0+\mathsf G_0\mathsf T\mathsf G_0.
\label{eq:feqo-maxwell-t}
\end{equation}
因此散射场为
\begin{equation}
\vb{E}_{\rm sc}\equiv\vb{E}-\vb{E}_{\rm inc}
=\mathsf G_0\mathsf T\vb{E}_{\rm inc}.
\label{eq:feqo-maxwell-scattered-field}
\end{equation}
对非磁性局域介电结构，背景与真实介电张量分别为
$\boldsymbol\varepsilon_h$、$\boldsymbol\varepsilon$，有
$\mathsf V=(\omega^2/c^2)(\boldsymbol\varepsilon-\boldsymbol\varepsilon_h)$。
介电率允许为复数；推迟边界条件与上述可逆条件仍须同时满足。
\end{proposition}
\begin{proof}
真实源方程为 $(\mathsf L_0-\mathsf V)\vb{E}=f$。
左乘背景逆，得到 $\vb{E}-\mathsf G_0\mathsf V\vb{E}=\mathsf G_0f$，
即\eqnrefs{eq:feqo-maxwell-lippmann}。
设 $\mathsf A=\mathsf G_0\mathsf V$，由假定的有界逆可求
$\vb{E}=(I-\mathsf A)^{-1}\vb{E}_{\rm inc}$。
这个场也属于原微分算符的定义域：因为
$\vb{E}=\mathsf G_0f+\mathsf G_0\mathsf V\vb{E}$，
右侧两项均在 $D(\mathsf L_0)$ 中。

逆算符恒等式 $(I-\mathsf A)^{-1}-I=\mathsf A(I-\mathsf A)^{-1}$ 给出
\[
\mathsf G=(I-\mathsf G_0\mathsf V)^{-1}\mathsf G_0
=\mathsf G_0+\mathsf G_0\mathsf V
(I-\mathsf G_0\mathsf V)^{-1}\mathsf G_0.
\]
定义 $\mathsf T=\mathsf V(I-\mathsf G_0\mathsf V)^{-1}$ 后得到真实逆的表示；
对同一个背景场作用，则得到\eqnrefs{eq:feqo-maxwell-scattered-field}。
非磁背景与真实算符的旋度项相同，二者相减只留下
$(\omega^2/c^2)(\boldsymbol\varepsilon-\boldsymbol\varepsilon_h)$，所以介电对比的符号为正。
\end{proof}

无限空间的出射实频 Green 算符通常不是普通 $L^2$ 上的有界逆。
上述表示应先在适当的加权场与源空间、或有限对比区域的体积分方程中建立，
再取实际存在的 $\omega+\ii\eta$、$\eta\downarrow0$ 边界值。
平面波照明在颗粒区域内也按这一体积分表示处理。
若目标频率出现逆的极点，或材料与空间选择不满足有界性及可逆条件，
不能继续直接使用常规有界逆；必须保留有限虚部或采用相应的谱与辐射边界分析。

\begin{exbox}[Mie 与 Rayleigh 极限]
\eqnrange{eq:feqo-maxwell-t}{eq:feqo-maxwell-scattered-field} 把满足所给条件的线性结构散射写成 Maxwell $T$ 算符。现在把几何专门化为半径 $a$ 的均匀各向同性球：球内相对介电率为 $\varepsilon_s(\omega)$，外部宿主为无损、非磁均匀介质 $\varepsilon_h>0$。求该球的电磁多极散射系数，并在小尺寸极限计算电偶极响应和散射、吸收截面。

\textbf{解。} 宿主波数和尺寸参数定义为
\begin{equation*}
\kappa_h\equiv \sqrt{\varepsilon_h}\,\frac{\omega}{c},
\qquad
x\equiv \kappa_h a,
\qquad
m\equiv\sqrt{\frac{\varepsilon_s}{\varepsilon_h}}.
\end{equation*}
$\kappa_h$ 的单位为 ${\rm m^{-1}}$，而 $x$ 与 $m$ 无量纲。被动材料取使场在无穷远满足出射条件的平方根分支。

球面对任意空间转动不变，因此 $\mathsf T_{\rm EM}$ 与总角动量对易。选择矢量球面波作为基以后，不同角动量 $\ell$、磁量子数与 TE/TM 偏振不再相互混合；求解问题退化为每个 $\ell$ 通道在 $r=a$ 处匹配切向 $\vb{E}$ 与 $\vb{H}$。定义 Riccati--Bessel 函数
\begin{equation*}
\psi_\ell(z)\equiv z j_\ell(z),
\qquad
\xi_\ell(z)\equiv z h_\ell^{(1)}(z),
\end{equation*}
并用撇号表示对完整自变量求导。对非磁球，两个无量纲 Mie 通道系数为
\begin{align}
 a_\ell
 &=
 \frac{m\psi_\ell(mx)\psi_\ell'(x)
 -\psi_\ell(x)\psi_\ell'(mx)}
 {m\psi_\ell(mx)\xi_\ell'(x)
 -\xi_\ell(x)\psi_\ell'(mx)},
 \label{eq:mie-electric-coefficient-dp75}\\
 b_\ell
 &=
 \frac{\psi_\ell(mx)\psi_\ell'(x)
 -m\psi_\ell(x)\psi_\ell'(mx)}
 {\psi_\ell(mx)\xi_\ell'(x)
 -m\xi_\ell(x)\psi_\ell'(mx)}.
 \label{eq:mie-magnetic-coefficient-dp75}
\end{align}
$a_\ell$ 是电型多极通道，$b_\ell$ 是磁型多极通道；二者正是同一 Maxwell $T$ 算符在球对称多极基中的对角矩阵元。

对平面波入射，正交多极通道的总散射截面和消光截面分别为
\begin{align}
 {
 C_{\rm sca}
 =\frac{2\pi}{\kappa_h^2}
 \sum_{\ell=1}^{\infty}(2\ell+1)
 \left(|a_\ell|^2+|b_\ell|^2\right),}
 \label{eq:mie-scattering-cross-section-dp75}\\
 {
 C_{\rm ext}
 =\frac{2\pi}{\kappa_h^2}
 \sum_{\ell=1}^{\infty}(2\ell+1)
 \operatorname{Re}(a_\ell+b_\ell),}
 \label{eq:mie-extinction-cross-section-dp75}
\end{align}
并由能量守恒定义
\begin{equation*}
C_{\rm abs}\equiv C_{\rm ext}-C_{\rm sca}\ge0
\end{equation*}
用于被动球。三个截面单位均为 ${\rm m^2}$。\eqnrefs{eq:mie-extinction-cross-section-dp75} 是球几何中的光学定理；当材料无损时 $C_{\rm abs}=0$，因此消光恰等于散射。

只有在外部尺寸参数和球内相位延迟都小，
\begin{equation}
{x\ll1,\qquad |m|x\ll1,}
\label{eq:rayleigh-control-dp75}
\end{equation}
才可以进一步进入 Rayleigh 长波极限。采用全文 $\ee^{-\ii\omega t}$ 约定，$\ell=1$ 通道的小参数展开为
\begin{align*}
\psi_1(z)&=\frac{z^2}{3}+O(z^4),
&\psi_1'(z)&=\frac{2z}{3}+O(z^3),\\
\xi_1(z)&=-\frac{\ii}{z}+O(z),
&\xi_1'(z)&=\frac{\ii}{z^2}+O(1).
\end{align*}
代入\eqnrefs{eq:mie-electric-coefficient-dp75} 得
\begin{align*}
\text{分子}
&=\frac{2m(m^2-1)}{9}x^3+O(x^5),\\
\text{分母}
&=\frac{\ii m(m^2+2)}{3}+O(x^2),
\end{align*}
所以
\begin{equation}
{
 a_1
 =-\frac{2\ii}{3}
 \frac{\varepsilon_s-\varepsilon_h}
 {\varepsilon_s+2\varepsilon_h}
 x^3+O(x^5),
 \qquad
 b_1=O(x^5).}
\label{eq:rayleigh-mie-a1-dp75}
\end{equation}
更高电型多极满足 $a_\ell=O(x^{2\ell+1})$，因此电偶极是最低阶非零辐射通道。把\eqnrefs{eq:rayleigh-mie-a1-dp75} 代入\eqnrefs{eq:mie-scattering-cross-section-dp75}，得到
\begin{equation}
{
 C_{\rm sca}^{\rm Rayleigh}
 =\frac{8\pi}{3}\kappa_h^4a^6
 \left|
 \frac{\varepsilon_s-\varepsilon_h}
 {\varepsilon_s+2\varepsilon_h}
 \right|^2
 \left[1+O(x^2)\right].}
\label{eq:rayleigh-scattering-cross-section-dp75}
\end{equation}
若宿主在该频段近似无色散，$\kappa_h=\sqrt{\varepsilon_h}\,\omega/c$，于是\eqnrefs{eq:rayleigh-scattering-cross-section-dp75} 的最低阶项严格含 $\omega^4a^6$；Rayleigh 的蓝光强散射标度由 $a_1$ 的 $x^3$ 长波展开直接产生。

同一结果可以从准静态边值问题核对。令外加均匀场沿 $z$ 方向，球外和球内电势分别写成
\begin{equation*}
\Phi_{\rm out}
=-E_0r\cos\theta+\frac{A\cos\theta}{r^2},
\qquad
\Phi_{\rm in}=-BE_0r\cos\theta.
\end{equation*}
在 $r=a$ 连续匹配 $\Phi$ 与法向位移 $\varepsilon_0\varepsilon\,\partial_r\Phi$，两条线性方程给
\begin{equation*}
A=a^3E_0
\frac{\varepsilon_s-\varepsilon_h}
{\varepsilon_s+2\varepsilon_h}.
\end{equation*}
球外偶极电势为 $p\cos\theta/(4\pi\varepsilon_0\varepsilon_h r^2)$，故电偶极极化率
\begin{equation}
{
\alpha_E^{(0)}
=4\pi\varepsilon_0\varepsilon_h a^3
\frac{\varepsilon_s-\varepsilon_h}
{\varepsilon_s+2\varepsilon_h},
\qquad
\vb{p}=\alpha_E^{(0)}\vb{E}_0.}
\label{eq:rayleigh-static-polarizability-dp75}
\end{equation}
于是\eqnrefs{eq:rayleigh-scattering-cross-section-dp75} 也可写成
\begin{equation*}
C_{\rm sca}^{\rm Rayleigh}
=\frac{\kappa_h^4}
{6\pi\varepsilon_0^2\varepsilon_h^2}
\left|\alpha_E^{(0)}\right|^2
\left[1+O(x^2)\right].
\end{equation*}
对有吸收的球，$\operatorname{Im}\varepsilon_s>0$ 还产生材料吸收。若要在无损球中同时保持光学定理，不能只把静态 $\alpha_E^{(0)}$ 代入消光公式，因为辐射反作用在 $O(x^3)$ 修正极化率的虚部；等价地，应保留 Mie 系数的更高阶项或使用满足
\begin{equation*}
\operatorname{Im}\!\left[(\alpha_E)^{-1}\right]
=-\frac{\kappa_h^3}{6\pi\varepsilon_0\varepsilon_h}
\end{equation*}
的辐射修正极化率。若 $x$ 或 $|m|x$ 不再小，则电偶极近似失效，应回到完整 Mie 级数；若几何不再是球体，则应回到\eqnrange{eq:feqo-maxwell-lippmann}{eq:feqo-maxwell-scattered-field} 的一般 Maxwell 散射问题，并重新检查相应空间、边界值和可逆条件。
\end{exbox}

冲击参数依赖的指数因子表明，大波矢近场分量会被电子轨迹到界面的距离直接空间滤波。

\subsection{均匀介质的精确发射运动学与经典极限}

均匀透明介质给出检验一般相位匹配与 Green 张量辐射公式的最简单基准。先不忽略电子反冲。设介质模式色散为 $\kappa=n(\omega)\omega/c$，初态电子能量、动量为 $E_i,\vb{p}_i$，发射一个能量为 $\hbar\omega$、波矢为 $\vb{\kappa}$ 的模式量子后，末态满足
\begin{equation}
{
E_f=E_i-\hbar\omega,
\qquad
\vb{p}_f=\vb{p}_i-\hbar\vb{\kappa},
\qquad
E_j^2=m_{\mathrm e}^2c^4+c^2p_j^2.}
\label{eq:cherenkov_exact_mass_shell}
\end{equation}
这三个条件分别表达能量守恒、均匀介质中的平移动量守恒以及电子初末态的精确相对论质量壳；它们正是\chapref{chap:electron-foundations}的精确交换恒等式~\eqnrefs{eq:feqo-primer-recoil} 在发射支 $\sigma=-1$ 上的具体化。代入
\begin{equation*}
K^\mu=\left(\frac{\omega}{c},\vb{\kappa}\right),
\qquad
P_i^\mu=\left(\frac{E_i}{c},\vb{p}_i\right),
\end{equation*}
便有
\begin{align*}
0
&=-2P_i\!\cdot K+\hbar K^2\\
&=-2\left(\frac{E_i \omega}{c^2}-p_i \kappa\cos\theta\right)
+\hbar\left(\frac{\omega^2}{c^2}-\kappa^2\right).
\end{align*}
再使用介质传播支的色散 $\kappa=n(\omega)\omega/c$ 与 $p_ic^2=E_iv_i$，得到
\begin{equation}
{
\cos\theta
=\frac{c}{n(\omega)v_i}
\left[1+\frac{\hbar\omega}{2E_i}(n^2(\omega)-1)\right].}
\label{eq:cherenkov_exact_recoil_angle}
\end{equation}
因此这里真正新增的物理输入只有介质模式色散 $\kappa=n\omega/c$；电子反冲代数已经由一般质量壳一般表达式一次性解决。实传播通道要求右端绝对值不超过 1，所以有限反冲不仅改变发射角，也可能关闭某一频率的传播解。

\begin{exbox}[Cherenkov 辐射]
对具有实折射率 $n(\omega)$ 的均匀透明介质，求小反冲极限中的出射角、辐射阈值和单位长度能谱，并与保留反冲的同步条件比较。

\textbf{解。} 只有在
\begin{equation*}
\epsilon_{\rm rec}^{(C)}\equiv
\left|\frac{\hbar\omega}{2E_i}(n^2-1)\right|\ll1
\end{equation*}
时，才可忽略括号中的反冲修正。此时
\begin{equation}
{
\cos\theta_C=\frac1{\beta_{\rm rel}n(\omega)},}
\label{eq:classical_cherenkov_angle}
\end{equation}
并且只有 $\beta_{\rm rel}n>1$ 才有实传播角。这个零反冲的均匀介质发射现象通常称为 Cherenkov 辐射，$\theta_C$ 是其经典发射角。下图用 $x=\hbar\omega/E_i$ 比较精确结果与这一极限。
\begin{center}
\includegraphics[width=0.74\linewidth]{cherenkov_recoil_angle.png}\\[-2pt]
{\small 固定 $\beta_{\rm rel}=0.95$ 时由精确质量壳得到的发射角；实线采用\eqnrefs{eq:cherenkov_exact_recoil_angle}，虚线为 $x\to0$ 的经典 Cherenkov 极限。}
\end{center}

辐射强度由给定电流对 $\operatorname{Im}G^R$ 的投影确定。无损、各向同性、非磁介质中，傅里叶空间 Green 张量的横向在壳部分是
\[
\operatorname{Im}\mathbf G^R(\vb{\kappa},\omega)
=\pi(\mathbb I-\hat{\vb{\kappa}}\hat{\vb{\kappa}})
\delta(\kappa^2-n^2\omega^2/c^2).
\]
把它代入\eqnrefs{eq:homogeneous_green_loss_general}，再分别完成径向积分与角向 $\delta$ 函数积分，即得到 Frank--Tamm 谱。相关归一化与材料响应可对照~\cite{FrankTamm1937,RiveraKaminer2020Quasiparticles}。
\begin{equation}
{
\frac{\dd^2 W}{\dd L\,\dd\omega}
=\frac{e^2}{4\pi\varepsilon_0c^2}\,
\omega\left(1-\frac1{\beta_{\rm rel}^2n^2(\omega)}\right)
\Theta[\beta_{\rm rel}n(\omega)-1].}
\label{eq:frank_tamm_spectrum}
\end{equation}
经典谱因此由同一个在壳条件同时决定阈值与横向极化因子。有限反冲的量子谱不能通过在\eqnrefs{eq:frank_tamm_spectrum} 上经验替换角度获得，而必须回到精确跃迁电流或 $S$ 矩阵表达。

\medskip

单光子反冲可以在 $\hbar\omega/E_i\to0$ 时忽略；若介质同时均匀、透明且电子在所考察区间内保持恒速，量子跃迁电流退化为经典匀速电流。此时辐射阈值只取决于该电流的同步 Fourier 分量能否与介质传播质量壳相交。对沿 $z$ 方向运动的给定电流，时间与纵向坐标积分产生同步 $\delta$ 函数
\begin{equation*}
2\pi\delta(\omega-v\kappa_z)
=2\pi\delta(\omega-v\kappa\cos\theta).
\end{equation*}
均匀、无损、非磁介质的推迟 Green 张量在傅里叶空间的横向虚部为
\begin{equation*}
\operatorname{Im}G_{ij}(\vb{\kappa},\omega)
=\pi\left(\delta_{ij}-\widehat{\vb{\kappa}}_i\widehat{\vb{\kappa}}_j\right)
\delta(\kappa^2-\kappa_m^2),
\qquad
\kappa_m=\frac{n(\omega)\omega}{c}.
\end{equation*}
电流方向是 $\hat z$，所以横向投影算符收缩给
\begin{equation*}
\hat z_i(\delta_{ij}-\widehat{\vb{\kappa}}_i\widehat{\vb{\kappa}}_j)\hat z_j
=1-\cos^2\theta.
\end{equation*}
于是一般 Green 张量损耗积分中需要计算的几何部分是
\begin{equation*}
\int_0^\infty \kappa^2\dd\kappa\int\dd\Omega\,
(1-\cos^2\theta)
\delta(\kappa^2-\kappa_m^2)
\delta(\omega-v\kappa\cos\theta).
\end{equation*}
径向 $\delta$ 函数先给
\begin{equation*}
\int_0^\infty \kappa^2\dd\kappa\,\delta(\kappa^2-\kappa_m^2)
=\frac{\kappa_m}{2}.
\end{equation*}
再令 $u=\cos\theta$，完成方位角积分：
\begin{align*}
2\pi\int_{-1}^{1}\dd u\,(1-u^2)
\delta(\omega-v\kappa_m u)
&=\frac{2\pi}{v\kappa_m}
\left(1-\frac{\omega^2}{v^2\kappa_m^2}\right)
\Theta(v\kappa_m-\omega).
\end{align*}
阶跃函数来自 $\delta$ 函数根
$u_*=\omega/(v\kappa_m)$ 必须位于 $[-1,1]$，这正是 $\beta_{\rm rel}n>1$ 的阈值。把两个积分结果先相乘，几何积分本身为
\begin{equation*}
\frac{\kappa_m}{2}\,
\frac{2\pi}{v\kappa_m}
\left(1-\frac{\omega^2}{v^2\kappa_m^2}\right)
\Theta(v\kappa_m-\omega)
=\frac{\pi}{v}
\left(1-\frac1{\beta_{\rm rel}^2n^2}\right)
\Theta(\beta_{\rm rel}n-1).
\end{equation*}
\eqnrefs{eq:homogeneous_green_loss_general} 还带有
$2\mu_0e^2v^2/(\hbar v)$ 与傅里叶积分测度的 $(2\pi)^{-3}$，而上面横向 Green 张量虚部中的整体谱前因子还提供一个 $\pi$。因此
\begin{align*}
\frac{\dd^2 P}{\dd L\,\dd\omega}
&=\frac{2\mu_0e^2}{\hbar v}
\frac{v^2\pi}{(2\pi)^3}
\frac{\pi}{v}
\left(1-\frac1{\beta_{\rm rel}^2n^2}\right)
\Theta(\beta_{\rm rel}n-1),\\
&=\frac{\mu_0e^2}{4\pi\hbar}
\left(1-\frac1{\beta_{\rm rel}^2n^2}\right)
\Theta(\beta_{\rm rel}n-1).
\end{align*}
最后乘以单量子能量 $\hbar\omega$，并用
$\mu_0=1/(\varepsilon_0c^2)$，便得到\eqnrefs{eq:frank_tamm_spectrum}。因此 Cherenkov 阈值和 $1-1/(\beta_{\rm rel}^2n^2)$ 权重同时由同一个横向在壳谱投影算符给出：前者决定传播支路是否存在，后者是该支路的横向投影权重。
\end{exbox}

极性晶格、周期材料与混合准粒子都对应于同一因果材料算符加入新的内部自由度或空间周期性；其微观顶点由晶格正常模与电子耦合的共同基本对象构造。

\subsection{周期介质顶点}\label{chap:feqo-structured-media}
在周期材料中，电子既可能与电磁场作用，也可能激发晶格振动。\chapref{chap:general-interaction}已说明怎样量子化正常模；现在把模式坐标具体写成各原子或各原胞的位移。单个原胞内部的振动决定模式频率与耦合大小，各原胞之间的相位关系则决定能够交换的动量。把两部分相乘并对晶格求和，就能得到周期介质中的跃迁矩阵元。

声子并不是先验加入材料模型的新粒子，而是晶格小振动的正常模量子。经典起点仍是各离子偏离平衡位置的位移 $\vb{u}_{l\iota}$：把势能在平衡构型附近展开到二阶，海森矩阵给出不同离子位移之间的线性回复力；周期性再允许用 Bloch 波矢把不同晶胞的振动解耦；对每个 Bloch 波矢对角化剩余的晶胞内部矩阵，得到的每一个独立谐振子量子化以后才称为一个声子模。因而“动力学矩阵 $\to$ 正常模 $\to$ 谐振子 $\to$ 声子产生湮灭算符”是一条连续的经典到量子链。

对三维周期谐性晶格，质量矩阵与力常数的 Bloch 变换把\chapref{chap:general-interaction}的一般二次本征问题具体化为动力学矩阵问题。以 $\vb{\kappa}$ 标记第一布里渊区内的波矢，以 $\nu$ 标记正常模分支，以 $\iota$ 标记晶胞内部自由度；$\vb{e}_{\iota\nu}(\vb{\kappa})$ 是质量加权动力学算符的归一化本征矢。经典线性晶格首先由质量加权动力学矩阵确定正常模：
\begin{equation}
{
\sum_{\iota'}\mathsf D_{\iota\iota'}(\vb{\kappa})
\,\vb{e}_{\iota'\nu}(\vb{\kappa})
=\omega_{\vb{\kappa}\nu}^2
\vb{e}_{\iota\nu}(\vb{\kappa}).}
\label{eq:phonon_normal_mode_eigenproblem}
\end{equation}
对每个本征对 $(\vb{\kappa},\nu)$，把沿本征矢的质量加权位移幅记为经典正常坐标 $Q_{\vb{\kappa}\nu}$，其正则共轭动量记为 $P_{\vb{\kappa}\nu}$。选择本征矢的标准归一化后，二次晶格哈密顿量被对角化为
\begin{equation*}
H_{\rm ph}^{(2)}
=\frac12\sum_{\vb{\kappa}\nu}
\left(|P_{\vb{\kappa}\nu}|^2
+\omega_{\vb{\kappa}\nu}^2|Q_{\vb{\kappa}\nu}|^2\right).
\end{equation*}
因此每个 $(\vb{\kappa},\nu)$ 都是一个独立谐振子；量子化以后
\begin{equation*}
\hat Q_{\vb{\kappa}\nu}
=\sqrt{\frac{\hbar}{2\omega_{\vb{\kappa}\nu}}}
\left(\hat a_{\vb{\kappa}\nu}
+\hat a_{-\vb{\kappa}\nu}^{\dagger}\right).
\end{equation*}
把这个正常坐标沿质量加权本征矢变回各离子的实际位移，就得到

\begin{equation}
{
\hat{\vb{u}}_{l\iota}
=\frac1{\sqrt N}\sum_{\vb{\kappa}\nu}
\sqrt{\frac{\hbar}{2M_\iota\omega_{\vb{\kappa}\nu}}}
\left[
\vb{e}_{\iota\nu}(\vb{\kappa})\hat a_{\vb{\kappa}\nu}\ee^{\ii\vb{\kappa}\cdot\vb{R}_l}
+\vb{e}_{\iota\nu}^{*}(\vb{\kappa})\hat a_{\vb{\kappa}\nu}^{\dagger}\ee^{-\ii\vb{\kappa}\cdot\vb{R}_l}
\right].}
\label{eq:phonon_displacement_quantization}
\end{equation}
其中 $N$ 是晶胞数，$\vb{R}_l$ 是第 $l$ 个晶胞位置，$\omega_{\vb{\kappa}\nu}$ 是正常模频率。复共轭项保证位移算符厄米；对应自由声子哈密顿量是一组独立谐振子。\eqnrefs{eq:phonon_displacement_quantization} 因而只是把\chapref{chap:general-interaction}的一般模式坐标代入周期晶格的质量、本征矢与 Bloch 相位；它给出这一材料类中的单量子位移幅，而不改变前面已经建立的玻色对易关系和联合动力学。

\begin{theorem}[周期谐性格的正交正常模量子化]
设有限周期晶格质量为正，力常数实对称且平移不变，质量加权动力学矩阵在所保留振动子空间上正定。Bloch 本征矢正交归一并取相容共轭规范，则正常坐标的正则量子化给独立玻色模与\eqnrefs{eq:phonon_displacement_quantization}的厄米位移算符。整体平移等零频声学模须先单独分离，不能在 $\omega=0$ 套用 $1/\sqrt\omega$。
\end{theorem}
\begin{proof}
从平衡晶格的小振动开始。第 $l$ 个晶胞中第 $\iota$ 个离子的位移记为 $\vb{u}_{l\iota}$，质量为 $M_\iota$。谐性近似下的经典哈密顿量为
\[
H^{(2)}
=\frac12\sum_{l\iota}M_\iota|\dot{\vb{u}}_{l\iota}|^2
+\frac12\sum_{l\iota,l'\iota'}
\vb{u}_{l\iota}\cdot
\boldsymbol\Phi_{l\iota,l'\iota'}\vb{u}_{l'\iota'},
\]
其中 $\boldsymbol\Phi$ 是平衡位置处势能海森矩阵。周期性使力常数只依赖晶胞差，因此可以作 Bloch 变换。定义质量加权 Fourier 坐标
\[
\vb{q}_\iota(\vb{\kappa})
=\sqrt{\frac{M_\iota}{N}}
\sum_l\vb{u}_{l\iota}\ee^{-\ii\vb{\kappa}\cdot\vb{R}_l}.
\]
逆变换为
\[
\vb{u}_{l\iota}
=\frac1{\sqrt{NM_\iota}}
\sum_{\vb{\kappa}}
\vb{q}_\iota(\vb{\kappa})
\ee^{\ii\vb{\kappa}\cdot\vb{R}_l}.
\]
对固定 $\vb{\kappa}$，动力学矩阵本征问题
\[
\sum_{\iota'}\mathsf D_{\iota\iota'}(\vb{\kappa})
\vb{e}_{\iota'\nu}(\vb{\kappa})
=\omega_{\vb{\kappa}\nu}^2
\vb{e}_{\iota\nu}(\vb{\kappa})
\]
给出一组正交本征矢。采用
\[
\sum_\iota
\vb{e}_{\iota\nu}^*(\vb{\kappa})\cdot
\vb{e}_{\iota\nu'}(\vb{\kappa})
=\delta_{\nu\nu'}
\]
的归一化，把质量加权位移展开为
\[
\vb{q}_\iota(\vb{\kappa})
=\sum_\nu\vb{e}_{\iota\nu}(\vb{\kappa})
Q_{\vb{\kappa}\nu}.
\]
代回逆 Bloch 变换：
\[
\vb{u}_{l\iota}
=\frac1{\sqrt{NM_\iota}}
\sum_{\vb{\kappa}\nu}
\vb{e}_{\iota\nu}(\vb{\kappa})
Q_{\vb{\kappa}\nu}
\ee^{\ii\vb{\kappa}\cdot\vb{R}_l}.
\]
由于真实空间位移是实变量，$Q_{-\vb{\kappa},\nu}=Q_{\vb{\kappa}\nu}^*$。对实力常数有 $\mathsf D(-\vb{\kappa})=\mathsf D^*(\vb{\kappa})$；在相容的本征矢规范下可取
\[
\vb{e}_{\iota\nu}(-\vb{\kappa})
=\vb{e}_{\iota\nu}^*(\vb{\kappa}).
\]

利用本征矢正交性，二次哈密顿量对角化为独立正常坐标：
\[
H^{(2)}
=\frac12\sum_{\vb{\kappa}\nu}
\left(P_{\vb{\kappa}\nu}P_{-\vb{\kappa}\nu}
+\omega_{\vb{\kappa}\nu}^2
Q_{\vb{\kappa}\nu}Q_{-\vb{\kappa}\nu}\right).
\]
量子化时令
\[
[\hat Q_{\vb{\kappa}\nu},
\hat P_{\vb{\kappa}'\nu'}]
=\ii\hbar\delta_{\nu\nu'}
\delta_{\vb{\kappa},-\vb{\kappa}'},
\]
并定义
\begin{align*}
\hat Q_{\vb{\kappa}\nu}
&=\sqrt{\frac{\hbar}{2\omega_{\vb{\kappa}\nu}}}
\left(\hat a_{\vb{\kappa}\nu}
+\hat a_{-\vb{\kappa}\nu}^\dagger\right),\\
\hat P_{\vb{\kappa}\nu}
&=-\ii\sqrt{\frac{\hbar\omega_{\vb{\kappa}\nu}}2}
\left(\hat a_{\vb{\kappa}\nu}
-\hat a_{-\vb{\kappa}\nu}^\dagger\right).
\end{align*}
由 $[a_{\vb{\kappa}\nu},a_{\vb{\kappa}'\nu'}^\dagger]
=\delta_{\vb{\kappa}\vb{\kappa}'}\delta_{\nu\nu'}$ 可直接验证上面的正则对易关系，并得到
\[
\hat H_{\rm ph}
=\sum_{\vb{\kappa}\nu}
\hbar\omega_{\vb{\kappa}\nu}
\left(\hat a_{\vb{\kappa}\nu}^\dagger
\hat a_{\vb{\kappa}\nu}+\frac12\right).
\]

最后把 $\hat Q$ 代回真实位移。第一部分为
\[
\frac1{\sqrt N}\sum_{\vb{\kappa}\nu}
\sqrt{\frac{\hbar}{2M_\iota\omega_{\vb{\kappa}\nu}}}
\vb{e}_{\iota\nu}(\vb{\kappa})
\hat a_{\vb{\kappa}\nu}
\ee^{\ii\vb{\kappa}\cdot\vb{R}_l}.
\]
第二部分对求和变量作 $\vb{\kappa}\mapsto-\vb{\kappa}$，并使用
$\omega_{-\vb{\kappa}\nu}=\omega_{\vb{\kappa}\nu}$ 与本征矢共轭关系，得到
\[
\frac1{\sqrt N}\sum_{\vb{\kappa}\nu}
\sqrt{\frac{\hbar}{2M_\iota\omega_{\vb{\kappa}\nu}}}
\vb{e}_{\iota\nu}^*(\vb{\kappa})
\hat a_{\vb{\kappa}\nu}^\dagger
\ee^{-\ii\vb{\kappa}\cdot\vb{R}_l}.
\]
两部分相加即为正文\eqnrefs{eq:phonon_displacement_quantization}。
\end{proof}


每个晶胞含两个离子的一维最近邻链可作为这个一般本征问题的二维内部空间实例；其两条本征支分别在长波极限表现为声学支和有限频率光学支。一般声子量子化仍由\eqnrefs{eq:phonon_displacement_quantization} 给出；最近邻链的力常数与质量具体指定动力学矩阵。

若给定的微观单电子哈密顿量 $\hat h[\{\vb{u}_{l\iota}\}]$ 随离子位置变化，则在平衡构型附近作 Taylor 展开，其一阶项包含该模型内全部线性位移耦合
\begin{equation*}
\delta\hat h
=\sum_{l\iota}\hat{\vb{u}}_{l\iota}\cdot
\left.\frac{\partial\hat h}{\partial\vb{u}_{l\iota}}\right|_0
=\sum_{\vb{\kappa}\nu}
\left(\hat V_{\vb{\kappa}\nu}\hat a_{\vb{\kappa}\nu}
+\hat V_{\vb{\kappa}\nu}^{\dagger}\hat a_{\vb{\kappa}\nu}^{\dagger}\right).
\end{equation*}
因此，在这一类由离子坐标参数化的单电子模型中，线性微观耦合的基本对象是 $\partial\hat h/\partial\vb{u}_{l\iota}$。形变势、Fr\"ohlich 耦合或外部自由电子对振动极化场的耦合，都是在不同电子 Hilbert 空间和不同长波近似下对这个耦合算符取矩阵元。

极性晶体中的光学声子还会产生宏观电极化。先区分两个关键的长波模：TO模的离子相对位移在 $\vb{\kappa}\to0$ 时垂直于传播方向，LO模的位移则平行于传播方向。极性晶体中的纵向位移会建立宏观 Coulomb 电场，因此长波 LO 与 TO 频率一般分裂；$\omega_{\rm TO}$ 是介电函数的共振极点，而无阻尼的 $\omega_{\rm LO}$ 对应纵向介电响应的零点。

令 $u$ 是一个红外活性 TO 正常坐标，$M_r$ 是该模的约化质量。宏观极化对离子相对位移的一阶导数定义 Born 有效电荷；在当前单模标量记号中写为 $Z^*e$，即单位位移产生的极化系数由 $Z^*e$ 给出。$n_{\rm ion}$ 是单位体积振子数，$\Gamma_{\rm ph}$ 是阻尼率。在线性响应中
\begin{equation}
{
M_r(\ddot u+\Gamma_{\rm ph}\dot u+\omega_{\rm TO}^2u)=Z^*eE,
\qquad
P_{\rm ion}=n_{\rm ion}Z^*eu.}
\label{eq:polar-optical-oscillator-eom-dp76}
\end{equation}
配合 $\vb{D}=\varepsilon_0\varepsilon_\infty\vb{E}+\vb{P}_{\rm ion}$ 得到单振子介电函数
\begin{equation}
{
\varepsilon(\omega)
=\varepsilon_\infty
\frac{\omega_{\rm LO}^2-\omega^2-\ii\Gamma_{\rm ph}\omega}
{\omega_{\rm TO}^2-\omega^2-\ii\Gamma_{\rm ph}\omega},
\qquad
\omega_{\rm LO}^2
=\omega_{\rm TO}^2+
\frac{n_{\rm ion}Z^{*2}e^2}{\varepsilon_0\varepsilon_\infty M_r}.}
\label{eq:polar_optical_dielectric_response}
\end{equation}
于是无阻尼静态极限给出 Lyddane--Sachs--Teller 关系
\begin{equation}
{
\frac{\varepsilon_s}{\varepsilon_\infty}
=\frac{\omega_{\rm LO}^2}{\omega_{\rm TO}^2}.}
\label{eq:lattice_lst_relation}
\end{equation}
它只描述材料响应的极点与零点；外部电子真正看到的耦合强度还要由模式场或 Green 张量与电子跃迁电流的投影决定。

\begin{exbox}[热声子吸收与发射]
取热平衡中的谐性晶格，设声子频率与电子--声子矩阵元已经由正常模展开给定。求电子吸收、发射一个声子的跃迁率，并比较两类过程的热占据权重。

\textbf{解。} 把\chapref{chap:electron-finite-time}\eqnrefs{eq:universal_boson_finite_time_emission} 的模式标签取为 $\lambda=(\vb{\kappa},\nu)$，并在长时间弱耗尽极限取 Fermi 黄金律。令
\begin{equation*}
 \mathcal M_{fi}^{(\vb{\kappa}\nu)}
 \equiv\langle f|\hat{\mathcal V}_{\vb{\kappa}\nu}|i\rangle
\end{equation*}
为去掉玻色升降矩阵元后的电子--模式耦合。热平衡平均占据数为 $N_{\vb{\kappa}\nu}$ 时，单个已指定电子末态的吸收与发射速率密度分别为
\begin{align*}
\Gamma_{\rm abs}^{i\to f}
&=\frac{2\pi}{\hbar}
\left|\mathcal M_{fi}^{(\vb{\kappa}\nu)}\right|^2
N_{\vb{\kappa}\nu}\,
\delta(E_f-E_i-\hbar\omega_{\vb{\kappa}\nu}),\\
\Gamma_{\rm em}^{i\to f}
&=\frac{2\pi}{\hbar}
\left|\mathcal M_{fi}^{(\vb{\kappa}\nu)}\right|^2
(N_{\vb{\kappa}\nu}+1)\,
\delta(E_f-E_i+\hbar\omega_{\vb{\kappa}\nu}).
\end{align*}
热平衡玻色占据数为
\begin{equation*}
N_{\vb{\kappa}\nu}
=\frac{1}{\exp(\beta\hbar\omega_{\vb{\kappa}\nu})-1},
\qquad \beta=(k_BT)^{-1}.
\end{equation*}
对同一模式、互为反过程的一对能量壳通道，除去共同电子矩阵元和态密度以后，发射与吸收的比值为
\begin{align*}
\frac{\Gamma_{\rm em}}{\Gamma_{\rm abs}}
&=\frac{N_{\vb{\kappa}\nu}+1}{N_{\vb{\kappa}\nu}}\\
&=\frac{\ee^{\beta\hbar\omega}/(\ee^{\beta\hbar\omega}-1)}{1/(\ee^{\beta\hbar\omega}-1)}\\
&=\ee^{\beta\hbar\omega_{\vb{\kappa}\nu}}.
\end{align*}
因此低温 $k_BT\ll\hbar\omega$ 时吸收被 Boltzmann 因子强烈压低，而自发发射仍保留 $+1$；高温 $k_BT\gg\hbar\omega$ 时 $N\gg1$，两者趋近相等。这里的 $+1$ 正是通用产生算符矩阵元 $|\langle N+1|a^\dagger|N\rangle|^2=N+1$ 的声子实例。有限作用时间、强耗尽或不可分辨连续谱时，应回到\chapref{chap:electron-finite-time}的有限时间传播与\chapref{chap:general-interaction}的多模振幅层级。
\end{exbox}

周期性还改变动量选择。对沿 $z$ 的一维周期 $\Lambda$，倒格矢为 $G_m=2\pi m/\Lambda$。这里必须区分周期方向与被结构破坏的横向平移对称性：离散 $z$ 平移只固定纵向 Bloch 准动量模倒格矢守恒，并不能单独固定电子末态的全部三维动量。若发射一个真空光子
$\vb{\kappa}=(\omega/c)\hat{\vb{n}}$，一般的能量壳与纵向准动量约束为
\begin{equation}
{
E_f=E_i-\hbar\omega,
\qquad
p_{f,z}=p_i-\hbar G_m-\hbar\kappa_z,
\qquad
E_f^2=m_{\mathrm e}^2c^4+c^2\!\left(p_{f,z}^2+p_{f,\perp}^2\right).}
\label{eq:periodic_photon_exact_conservation}
\end{equation}
结构若在某个横向方向不具连续平移对称性，就可以吸收该方向的连续动量，因此 $p_{f,\perp}$ 必须由具体边界条件、Green 张量或末态通道另行决定。仅由\eqnrefs{eq:periodic_photon_exact_conservation} 一般不存在唯一的有限反冲 $\omega(\theta)$ 闭式。

为了说明一个可封闭求解的特例，再考虑理想化通道：除周期方向的 $\hbar G_m$ 外，结构不交换横向动量，于是
\begin{equation*}
\vb{p}_{f,\perp}=-\hbar\vb{\kappa}_\perp .
\end{equation*}
在这个附加通道条件下，三维末态可写成
\begin{equation*}
\vb{p}_f=(p_i-\hbar G_m)\hat{\vb{z}}-\hbar\vb{\kappa}.
\end{equation*}
把传播光子和倒格动量合成该指定通道中的有效正频交换四矢量
\begin{equation*}
\mathcal K_m^\mu
\equiv
\left(
\frac{\omega}{c},
\vb{\kappa}+G_m\hat{\vb{z}}
\right),
\qquad
P_f=P_i-\hbar\mathcal K_m .
\end{equation*}
其中 $\mathcal K_m^\mu$ 表示该指定通道中的总四动量交换；它的空间部分同时含传播光子的动量与晶格提供的倒格动量。由发射支的统一质量壳恒等式
$-2P_i\!\cdot\mathcal K_m+\hbar\mathcal K_m^2=0$，并代入
$\vb{\kappa}=(\omega/c)\hat{\vb{n}}$、
$\hat{\vb{z}}\!\cdot\hat{\vb{n}}=\cos\theta$，逐项得到
\begin{align*}
0
&=-2\left[
\frac{E_i \omega}{c^2}
-p_i\left(G_m+\frac{\omega}{c}\cos\theta\right)
\right]\\
&\quad
+\hbar\left[
-G_m^2-2G_m\frac{\omega}{c}\cos\theta
\right],
\end{align*}
其中真空光子的 $\omega^2/c^2-\kappa^2$ 已精确相消。整理为
\begin{equation*}
\omega\left[E_i-c(p_i-\hbar G_m)\cos\theta\right]
=c^2G_m\left(p_i-\frac{\hbar G_m}{2}\right),
\end{equation*}
所以
\begin{equation}
{
\omega
=\frac{c^2G_m\left(p_i-\hbar G_m/2\right)}
{E_i-c\left(p_i-\hbar G_m\right)\cos\theta}.}
\label{eq:smith_purcell_exact_recoil_dispersion}
\end{equation}
这里 $\theta$ 是光子方向与电子初始运动方向的夹角。与均匀介质相比，周期结构新增的物理输入不是另一套质量壳代数，而是离散平移对称性允许的 $G_m$ 以及指定横向通道条件。\eqnrefs{eq:smith_purcell_exact_recoil_dispersion} 因而只属于这个可封闭特例，不是任意表面光栅几何的普适有限反冲色散。

\begin{exbox}[Smith--Purcell 色散]
取沿电子运动方向周期为 $\Lambda$ 的结构，电子以速度 $v_i$ 匀速通过，末态光子射向真空。求第 $m$ 个倒格分量的小反冲同步色散，并写出适用条件。

\textbf{解。} 真正具有更广适用性的 Smith--Purcell 主关系来自小反冲下的纵向同步。若显著末态满足
\begin{equation*}
\epsilon_{\rm rec}\equiv
\frac{\max\{\hbar|G_m|,\,\hbar\omega/c,\,|p_{f,\perp}|\}}{p_i}\ll1,
\end{equation*}
则电子能量差的一阶展开只保留初始速度方向的动量变化：
\begin{align*}
E_i-E_f
&=v_i(p_i-p_{f,z})+\order(\epsilon_{\rm rec}^2E_i),\\
&=\hbar v_i(G_m+\kappa_z)+\order(\epsilon_{\rm rec}^2E_i).
\end{align*}
与 $E_i-E_f=\hbar\omega$ 以及 $\kappa_z=(\omega/c)\cos\theta$ 联立，得到
\begin{equation*}
\omega\simeq\frac{v_iG_m}{1-\beta_{\rm rel}\cos\theta}.
\end{equation*}
再用 $G_m=2\pi m/\Lambda$ 与 $\lambda_{\rm rad}=2\pi c/\omega$，
\begin{equation}
{
\lambda_{\rm rad}
=\frac{\Lambda}{m}
\left(\frac1{\beta_{\rm rel}}-\cos\theta\right).}
\label{eq:smith_purcell_wavelength}\end{equation}
这一无反冲周期辐射关系通常称为 Smith--Purcell 色散。它由能量守恒与周期方向的 Bloch 准动量匹配共同确定：电子同步线被倒格矢平移后与真空光锥相交。这里 $G_m$ 表示晶格离散平移对称性允许的准动量差；传播光子的四动量仍由 $(\hbar\omega/c,\hbar\vb{\kappa})$ 给出。
\end{exbox}

\begin{proposition}[有限周期驱动的同步窗口]
对一个 Bloch 场分量 $u_m\ee^{\ii(k_z+G_m)z}$，在长度 $L$ 的均匀窗口内，轨迹振幅因子为
\begin{equation}
\int_{-L/2}^{L/2}\ee^{\ii\Delta k z}\dd z
=L\sinc(\Delta kL/2),
\qquad\Delta k=k_z+G_m-\omega/v.
\label{eq:feqo-periodic-phase-matching}
\end{equation}
其强度主瓣第一零点为 $|\Delta k|=2\pi/L$，全半高宽为 $5.56623\ldots/L$。同步是允许条件；若模式重叠 $u_m=0$，该通道仍无激发。
\end{proposition}
\begin{proof}
$\Delta k\ne0$ 时直接积分指数得 $2\sin(\Delta kL/2)/\Delta k$；零点处由连续性得到 $L$。强度归一化后为 $\sinc^2(\Delta kL/2)$，其第一半高根 $x_{1/2}=1.391557\ldots$ 给出全宽 $4x_{1/2}/L$。
\end{proof}

\begin{workedexample}[垂直出射与有限角度宽度]
设电子速度为 $0.70c$，目标真空波长为 $800\,\mathrm{nm}$，取一级周期散射和垂直出射。求所需周期，并计算长度为 $L$ 的均匀相互作用窗产生的角度宽度。

\textbf{解。} 对真空出射光 $k_z=(\omega/c)\cos\theta$，小反冲下的同步式恢复\eqnrefs{eq:smith_purcell_wavelength}。在 $\theta=\pi/2$ 时，周期 $\Lambda=m\lambda_{\rm rad}v/c$。若光波长为 $800\,\mathrm{nm}$、$v=0.70c$、$m=1$，周期为 $560\,\mathrm{nm}$。

固定频率时，中心到第一零点的方向余弦差为 $|\Delta\cos\theta|=\lambda_{\rm rad}/L$；强度全半高宽为 $5.56623\lambda_{\rm rad}/(2\pi L)$。这些是理想均匀窗口的宽度，还需乘模式耦合和收集效率。

无限匀速亚光速电子的同步线位于真空光锥外。周期倒格矢是把它搬入传播区的一种机制；有限路径、加速与边缘也会改变源谱或平移对称性，应由相应边界问题计算。
\end{workedexample}

\begin{figure}[htbp]
\centering
\includegraphics[width=.72\linewidth]{light_cone_synchronism.pdf}
\caption{亚光速电子同步线与真空光锥。无限平移对称几何中，两者除原点外不相交；周期结构可通过倒格动量平移源谱，有限路径及加速则需使用对应的源与边界条件。}
\label{fig:light-cone-synchronism}
\end{figure}


有限周期数 $N$ 不改变中心色散，却决定相干线形。若相邻周期的相位失配为
\begin{equation*}
\Delta\varphi_{\rm cell}=\Lambda\left(\frac{\omega}{v}-\kappa_z-G_m\right),
\end{equation*}
则阵列因子为
\begin{equation*}
\left|\sum_{n=0}^{N-1}\ee^{\ii n\Delta\varphi_{\rm cell}}\right|^2
=\left[\frac{\sin(N\Delta\varphi_{\rm cell}/2)}{\sin(\Delta\varphi_{\rm cell}/2)}\right]^2.
\end{equation*}
理想相干极限中峰值按 $N^2$ 增长、主峰宽度按 $1/N$ 缩窄；位置无序、损耗或周期路径可分辨信息会压低不同周期之间的交叉项。

最后必须区分同一声子模面对两类电子时的矩阵元。材料内部 Bloch 电子可以在长波近似下用形变势或 Fr\"ohlich 顶点参数化，例如各向同性声学形变势给出
\begin{equation*}
M_{\vb{\kappa}}=D\sqrt{\frac{\hbar|\vb{\kappa}|}{2\rho_MVs}},
\end{equation*}
而极性 LO 模的长程 Coulomb 顶点满足
\begin{equation*}
|M_{\vb{\kappa}}^{\rm F}|^2
=\frac{e^2\hbar\omega_{\rm LO}}{2\varepsilon_0V|\vb{\kappa}|^2}
\left(\frac1{\varepsilon_\infty}-\frac1{\varepsilon_s}\right).
\end{equation*}
这些参数描述的是晶体能带态。样品外的自由电子则必须使用真空侧模式势或 Green 张量与自由电子波函数的矩阵元，典型的一量子模式顶点为
\begin{equation*}
\hbar\mathcal G_{\vb{p}'s',\vb{p}s}^{(\lambda)}
=\int \dd^3 r\,\psi_{\vb{p}'s'}^\dagger
\left[e\phi_\lambda^{\rm out}-ec\,\boldsymbol\alpha\cdot\vb{A}_\lambda^{\rm out}\right]
\psi_{\vb{p}s},
\end{equation*}
必要时再加入非电磁短程势。两类问题共享同一晶格正常模，但不共享同一个电子矩阵元。

有限阵列的普适振幅首先是各单元实际局域贡献的相干和，
\begin{equation*}
\mathscr A_N^{\rm arr}
=\sum_{n=0}^{N-1}\mathscr A_n^{\rm loc}.
\end{equation*}
只有当端点不会显著重塑局域模式、各周期单元的局域耦合可视为相同，并且有限截断的主要作用只是增加逐单元传播相位时，才可因子化为
\begin{equation}
{
\mathscr A_N^{\rm arr}
=\mathscr A_1^{\rm cell}\sum_{n=0}^{N-1}\ee^{\ii n\varphi_{\rm cell}}.}
\label{eq:finite-periodic-array-amplitude-dp68}
\end{equation}
在这一受控的等价单元近似下，有限长度把无限结构的严格相位匹配条件展宽成有限阵列因子
\begin{equation}
{
\mathscr D_N(\varphi_{\rm cell})
=\frac{\sin^2(N\varphi_{\rm cell}/2)}
{\sin^2(\varphi_{\rm cell}/2)}.}
\label{eq:finite-periodic-array-factor-dp68}
\end{equation}
其相位匹配极限为 $\mathscr D_N(2\pi m)=N^2$，第一零点距离主峰 $O(1/N)$。若端点改变局域场、单元间存在不可忽略的多重散射，或每个单元的耦合幅度明显不同，就不能使用单一 $\mathscr A_1^{\rm cell}$；此时应回到逐单元振幅和，或直接求有限结构的 Maxwell Green 张量。\eqnrefs{eq:finite-periodic-array-factor-dp68} 因而是等价单元因子化后的几何结果，而不是任意有限周期结构的普适动力学定理。


\begin{derivation}[远区角通量积分与总辐射功率]

从远区辐射场
\[
\vb{E}_{\rm rad}
=
\frac{q_{\mathrm e}}{4\pi\varepsilon_0cR}
\frac{
\vb{n}\times[
(\vb{n}-\vb{\beta}_{\rm rel})
\times\dot{\vb{\beta}}_{\rm rel}]
}
{\mathfrak d_{\rm ret}^{\,3}},
\qquad
\mathfrak d_{\rm ret}
=
1-\vb{n}\cdot\vb{\beta}_{\rm rel}
\]
出发。远区场横向于传播方向，因此
\[
\vb{n}\cdot\vb{E}_{\rm rad}=0,
\qquad
\vb{B}_{\rm rad}
=
\frac1c\vb{n}\times\vb{E}_{\rm rad}.
\]
Poynting 向量的径向分量于是为
\[
\vb{n}\cdot\vb{S}_{\rm rad}
=
\frac1{\mu_0}
\vb{n}\cdot
(\vb{E}_{\rm rad}\times\vb{B}_{\rm rad})
=
\varepsilon_0c|\vb{E}_{\rm rad}|^2.
\]
乘球面元 \(R^2\dd\Omega\)，按\emph{观察时间} \(t\) 计的角能流为
\[
\frac{\dd P_{\rm obs}}{\dd\Omega}
=
\frac{q_{\mathrm e}^2}{16\pi^2\varepsilon_0c}
\frac{
\left|
\vb{n}\times[
(\vb{n}-\vb{\beta}_{\rm rel})
\times\dot{\vb{\beta}}_{\rm rel}]
\right|^2
}
{\mathfrak d_{\rm ret}^{\,6}}.
\]

这里还不能直接把它称为“源在单位时间内发射的功率”，因为不同观察方向对应的推迟时间间隔不同。由
\[
t-t_{\rm r}
=
\frac{R(t_{\rm r})}{c}
\]
求导，并用
\[
\frac{\dd R}{\dd t_{\rm r}}
=
-\vb{n}\cdot\vb{v}
\]
得到
\[
\frac{\dd t}{\dd t_{\rm r}}
=
1-\vb{n}\cdot\vb{\beta}_{\rm rel}
=
\mathfrak d_{\rm ret}.
\]
因此同一辐射能量满足
\[
\frac{\dd E}{\dd t_{\rm r}\dd\Omega}
=
\frac{\dd E}{\dd t\dd\Omega}
\frac{\dd t}{\dd t_{\rm r}},
\]
从而
\[
{
\frac{\dd P_{\rm emit}}{\dd\Omega}
=
\frac{q_{\mathrm e}^2}{16\pi^2\varepsilon_0c}
\frac{
\left|
\vb{n}\times[
(\vb{n}-\vb{\beta}_{\rm rel})
\times\dot{\vb{\beta}}_{\rm rel}]
\right|^2
}
{(1-\vb{n}\cdot\vb{\beta}_{\rm rel})^5}.
}
\]
这就是正文\eqnrefs{eq:lienard_angular_emission_power}。分母从六次幂降到五次幂，正是推迟时间雅可比因子的结果。

总辐射功率可由这个角分布直接积分得到。由于总功率只依赖 \(\vb{\beta}_{\rm rel}\) 与 \(\dot{\vb{\beta}}_{\rm rel}\) 的相对方向，可取
\[
\vb{\beta}_{\rm rel}
=
\beta_{\rm rel}\vb{e}_z,
\qquad
\dot{\vb{\beta}}_{\rm rel}
=
\dot\beta_\perp\vb{e}_x
+
\dot\beta_\parallel\vb{e}_z,
\]
并写
\[
u=\cos\theta,
\qquad
\dd\Omega=\dd\varphi\,\dd u.
\]
对方位角先积分。直接展开双重叉积并使用
\[
\int_0^{2\pi}\cos\varphi\,\dd\varphi=0,
\qquad
\int_0^{2\pi}\cos^2\varphi\,\dd\varphi=\pi
\]
可得
\begin{align*}
&\int_0^{2\pi}\dd\varphi\,
\left|
\vb{n}\times[
(\vb{n}-\vb{\beta}_{\rm rel})
\times\dot{\vb{\beta}}_{\rm rel}]
\right|^2\\
&\qquad=
\pi\Bigl[
2\dot\beta_\parallel^2(1-u^2)
+
\dot\beta_\perp^2
\bigl((1+\beta_{\rm rel}^2)(1+u^2)-4\beta_{\rm rel}u\bigr)
\Bigr].
\end{align*}
平行、垂直两部分因此分别需要两个一维积分：
\[
I_\parallel
=
2\pi
\int_{-1}^{1}
\frac{1-u^2}{(1-\beta_{\rm rel}u)^5}\,\dd u,
\]
\[
I_\perp
=
\pi
\int_{-1}^{1}
\frac{
(1+\beta_{\rm rel}^2)(1+u^2)-4\beta_{\rm rel}u
}
{(1-\beta_{\rm rel}u)^5}\,\dd u.
\]
这些是有理函数积分。令
\[
w=1-\beta_{\rm rel}u,
\qquad
u=\frac{1-w}{\beta_{\rm rel}},
\qquad
\dd u=-\frac{\dd w}{\beta_{\rm rel}},
\]
上下限变为 \(w=1+\beta_{\rm rel}\) 与
\(w=1-\beta_{\rm rel}\)。逐项积分 \(w^{-3},w^{-4},w^{-5}\) 后得到
\[
I_\parallel
=
\frac{8\pi}{3}
\frac1{(1-\beta_{\rm rel}^2)^3}
=
\frac{8\pi}{3}\gamma_{\rm L}^6,
\]
以及
\[
I_\perp
=
\frac{8\pi}{3}
\frac1{(1-\beta_{\rm rel}^2)^2}
=
\frac{8\pi}{3}\gamma_{\rm L}^4.
\]
代回角功率，
\begin{align*}
P_{\rm rad}
&=
\frac{q_{\mathrm e}^2}{16\pi^2\varepsilon_0c}
\left(
I_\parallel\dot\beta_\parallel^2
+
I_\perp\dot\beta_\perp^2
\right)\\
&=
\frac{q_{\mathrm e}^2}{6\pi\varepsilon_0c}
\left(
\gamma_{\rm L}^6\dot\beta_\parallel^2
+
\gamma_{\rm L}^4\dot\beta_\perp^2
\right).
\end{align*}
由于
\[
\dot{\vb{\beta}}_{\rm rel}
=
\frac{\vb{a}}{c},
\qquad
|\vb{\beta}_{\rm rel}\times\vb{a}|^2
=
\beta_{\rm rel}^2a_\perp^2,
\]
并且
\[
\gamma_{\rm L}^6
\left[
a^2-
|\vb{\beta}_{\rm rel}\times\vb{a}|^2
\right]
=
\gamma_{\rm L}^6a_\parallel^2
+
\gamma_{\rm L}^4a_\perp^2,
\]
最终得到
\[
{
P_{\rm rad}
=
\frac{q_{\mathrm e}^2\gamma_{\rm L}^6}
{6\pi\varepsilon_0c^3}
\left[
a^2-
|\vb{\beta}_{\rm rel}\times\vb{a}|^2
\right],
}
\]
即正文\eqnrefs{eq:lienard_total_power_from_angular_flux}。非相对论极限
\(\beta_{\rm rel}\to0\) 自动退化为 Larmor 功率
\[
P_{\rm rad}
\longrightarrow
\frac{q_{\mathrm e}^2a^2}{6\pi\varepsilon_0c^3}.
\]
\end{derivation}

\begin{derivation}[超相对论辐射的形成长度]
真空远场中 $\kappa_z=(\omega/c)\cos\theta$，而 $v=\beta_{\rm rel}c$。因此正文\eqnrefs{eq:formation_length_general_mismatch} 精确给出
\begin{equation*}
L_f
=\frac{c}{\omega\,[\beta_{\rm rel}^{-1}-\cos\theta]}.
\end{equation*}
令 $x=\gamma_{\rm L}^{-2}$，则
$\beta_{\rm rel}^{-1}=(1-x)^{-1/2}$。对函数
$f(x)=(1-x)^{-1/2}$ 在 $x=0$ 处作一阶 Taylor 展开。由于
\begin{equation*}
f''(u)=\frac{3}{4}(1-u)^{-5/2},
\end{equation*}
Lagrange 余项给出某个 $\xi\in(0,x)$ 使
\begin{equation*}
r_\beta(x)
=f(x)-1-\frac{x}{2}
=\frac{f''(\xi)}{2}x^2.
\end{equation*}
因 $f''(u)$ 在 $0\le u<x<1$ 上单调递增，
\begin{equation*}
0\le r_\beta(x)
\le\frac{3x^2}{8(1-x)^{5/2}}.
\end{equation*}

角度部分使用四阶 Taylor 公式。存在介于 $0$ 与 $\theta$ 之间的中间角 $\theta_*$，使
\begin{equation*}
\cos\theta
=1-\frac{\theta^2}{2}
+\frac{\theta^4}{24}\cos\theta_*.
\end{equation*}
所以
\begin{equation*}
r_\theta(\theta)
=1-\cos\theta-\frac{\theta^2}{2}
=-\frac{\theta^4}{24}\cos\theta_*,
\qquad
|r_\theta|\le\frac{\theta^4}{24}.
\end{equation*}
由此
\begin{align*}
\beta_{\rm rel}^{-1}-\cos\theta
&=\frac{x+\theta^2}{2}+r_\beta+r_\theta\\
&=\frac{x+\theta^2}{2}(1+\delta_f),
\end{align*}
其中
\begin{equation*}
\delta_f
=\frac{2(r_\beta+r_\theta)}{x+\theta^2}.
\end{equation*}
三角不等式给出
\begin{equation*}
|\delta_f|
\le
{
\epsilon_f
\equiv
\frac{
\dfrac{3x^2}{4(1-x)^{5/2}}
+\dfrac{\theta^4}{12}
}{x+\theta^2}}.
\end{equation*}
代回形成长度即得到正文\eqnrefs{eq:formation_length_ultrarelativistic}。若 $\epsilon_f<1$，则
\begin{align*}
\left|\frac{L_f-L_f^{(0)}}{L_f^{(0)}}\right|
&=\left|\frac1{1+\delta_f}-1\right|\\
&=\frac{|\delta_f|}{|1+\delta_f|}\\
&\le\frac{\epsilon_f}{1-\epsilon_f},
\end{align*}
其中
$L_f^{(0)}=2c/[\omega(x+\theta^2)]$。这给出了常用超相对论小角形成长度的定量适用条件。
\end{derivation}

\begin{derivation}[源损失中的吸收与辐射通道]
先证明 Green 表达式。对 $\ee^{-\ii\omega t}$ 的时间约定，频域 Maxwell 方程给
\[
\widetilde{\vb{E}}
=\ii\mu_0\omega\,\mathbf G^R\widetilde{\vb{J}}.
\]
定义复数
\[
Z_J
\equiv
\iint\dd^3r\,\dd^3r'\,
\widetilde{\vb{J}}^{*}(\vb{r})\cdot
\mathbf G^R(\vb{r},\vb{r}';\omega)
\cdot\widetilde{\vb{J}}(\vb{r}').
\]
则
\[
\int\dd^3r\,\widetilde{\vb{J}}^{*}\cdot\widetilde{\vb{E}}
=\ii\mu_0\omega Z_J.
\]
电子或外部源损失给场的平均功率取负的机械做功率，
\begin{align*}
P_{\rm loss}
&=-\frac12\Re(\ii\mu_0\omega Z_J)\\
&=\frac{\mu_0\omega}{2}\Im Z_J.
\end{align*}
而
\[
\Im Z_J
=\iint\dd^3r\,\dd^3r'\,
\widetilde{\vb{J}}^{*}(\vb{r})\cdot
\operatorname{Im}\mathbf G^R(\vb{r},\vb{r}';\omega)
\cdot\widetilde{\vb{J}}(\vb{r}'),
\]
得到正文\eqnrefs{eq:source_total_loss_green_power}。

局域能量守恒进一步把总损耗分解为材料吸收和辐射。取非磁介质
$\vb{B}=\mu_0\vb{H}$、
$\vb{D}=\varepsilon_0\boldsymbol\varepsilon_r\vb{E}$。频域 Maxwell 方程为
\[
\nabla\times\vb{E}=\ii\omega\mu_0\vb{H},
\qquad
\nabla\times\vb{H}=\vb{J}-\ii\omega\varepsilon_0
\boldsymbol\varepsilon_r\vb{E}.
\]
第二式取复共轭：
\[
\nabla\times\vb{H}^*
=\vb{J}^*+\ii\omega\varepsilon_0
\boldsymbol\varepsilon_r^\dagger\vb{E}^*.
\]
使用向量恒等式
\[
\nabla\cdot(\vb{E}\times\vb{H}^*)
=\vb{H}^*\cdot(\nabla\times\vb{E})
-\vb{E}\cdot(\nabla\times\vb{H}^*)
\]
得到
\begin{align*}
\nabla\cdot(\vb{E}\times\vb{H}^*)
&=\ii\omega\mu_0|\vb{H}|^2
-\vb{E}\cdot\vb{J}^*\\
&\quad
-\ii\omega\varepsilon_0
\vb{E}\cdot\boldsymbol\varepsilon_r^\dagger\vb{E}^*.
\end{align*}
对体积 $V$ 积分并取实部。第一项纯虚，不贡献平均能流；利用
\[
\Re\!\left[-\ii\,
\vb{E}\cdot\boldsymbol\varepsilon_r^\dagger\vb{E}^*\right]
=-\vb{E}^*\cdot\boldsymbol\varepsilon_I\cdot\vb{E},
\]
以及散度定理，得到
\begin{align*}
\frac12\Re\oint_{\partial V}\dd\vb{A}\cdot
(\vb{E}\times\vb{H}^*)
&=-\frac12\Re\int_V\dd^3r\,\vb{J}^*\cdot\vb{E}\\
&\quad
-\frac{\omega\varepsilon_0}{2}
\int_V\dd^3r\,
\vb{E}^*\cdot\boldsymbol\varepsilon_I\cdot\vb{E}.
\end{align*}
识别
\[
P_{\rm rad}
=\frac12\Re\oint\dd\vb{A}\cdot(\vb{E}\times\vb{H}^*),
\qquad
P_{\rm abs}
=\frac{\omega\varepsilon_0}{2}\int\dd^3r\,
\vb{E}^*\cdot\boldsymbol\varepsilon_I\cdot\vb{E},
\]
以及
$P_{\rm loss}=-(1/2)\Re\int\vb{J}^*\cdot\vb{E}$，就得到
\[
{P_{\rm loss}=P_{\rm abs}+P_{\rm rad}.}
\]
对被动介质 $\boldsymbol\varepsilon_I\succeq0$，故 $P_{\rm abs}\ge0$；采用出射边界条件时 $P_{\rm rad}\ge0$。
\end{derivation}

\begin{derivation}[各向异性导电薄层的模式极点]
先从面内平移不变情形开始。令
\[
\sigma_{ab}(\vb{\rho},\vb{\rho}';\omega)
=
\sigma_{ab}(\vb{\rho}-\vb{\rho}';\omega).
\]
采用二维 Fourier 变换
\[
f(\vb{\rho})
=
\int\frac{\dd^2q}{(2\pi)^2}
f(\vb{q})\ee^{+\ii\vb{q}\cdot\vb{\rho}},
\]
卷积定理把正文\eqnrefs{eq:general_2d_conductivity_kernel} 化成
\[
K_a(\vb{q},\omega)
=
\sum_b
\sigma_{ab}(\vb{q},\omega)
E_{\parallel,b}(\vb{q},0,\omega).
\]

对固定 $\vb{q}$，两侧体介质中的平面波都满足无源 Maxwell 方程。先看 TM 极化。取 $\vb{q}=q\hat{\vb{x}}$，切向电场为 $E_x$、磁场为 $H_y$。由
\[
\vb{H}
=
\frac1{\mu_0\omega}\vb{k}\times\vb{E}
\]
和
$\vb{k}\cdot\varepsilon_j\vb{E}=0$，对沿 $\pm z$ 传播的 TM 波得到
\[
H_y^{(\pm)}
=
\pm
\frac{\varepsilon_0\varepsilon_j\omega}{k_{zj}}
E_x^{(\pm)}
=
\pm Y_{p,j}E_x^{(\pm)}.
\]
TE 极化同理给
\[
H_x^{(\pm)}
=
\mp
\frac{k_{zj}}{\mu_0\omega}
E_y^{(\pm)}
=
\mp Y_{s,j}E_y^{(\pm)}.
\]
所以 $Y_{p,j}$、$Y_{s,j}$ 正是把切向电场映射到切向磁场的波导纳。

现在寻找无外部入射的束缚/共振表面模。取两侧都从界面向外衰减或满足出射条件的解，并令共同切向场幅为 $E_\parallel$。切向电场连续性已经满足。磁场跳跃条件
\[
\hat{\vb{z}}\times(\vb{H}_2-\vb{H}_1)
=
\vb{K}
=
\boldsymbol\sigma\vb{E}_\parallel
\]
则给出一个齐次线性方程。对各向同性二维层，TM 与 TE 不混合，因此分别得到
\[
\bigl(Y_{p,1}+Y_{p,2}+\sigma_L\bigr)E_p=0,
\]
\[
\bigl(Y_{s,1}+Y_{s,2}+\sigma_T\bigr)E_s=0.
\]
存在非零 $E_p$ 或 $E_s$ 的条件正是正文\eqnrefs{eq:2d_sheet_surface_mode_conditions}。

若
\[
\boldsymbol\sigma
=
\begin{pmatrix}
\sigma_{pp}&\sigma_{ps}\\
\sigma_{sp}&\sigma_{ss}
\end{pmatrix}
\]
在 TE/TM 基底中含非对角元，边界方程变成
\[
\left[
\begin{pmatrix}
Y_{p,1}+Y_{p,2}&0\\
0&Y_{s,1}+Y_{s,2}
\end{pmatrix}
+
\boldsymbol\sigma
\right]
\begin{pmatrix}E_p\\E_s\end{pmatrix}
=0.
\]
非零解存在当且仅当系数矩阵行列式为零，即正文\eqnrefs{eq:2d_sheet_anisotropic_pole_determinant}。

最后说明为什么二维空间色散对电子特别重要。电子在距界面有限冲击参数处的 Fourier 电流包含连续 $q$，且近场部分在
$q>\omega/c$ 时对应倏逝真空波。远场光只在接近光锥的横向通道取样，而 EELS 的 Green 核积分会覆盖这些大 $q$ 倏逝分量；它们通过上述边界矩阵直接读取 $\sigma_L(q,\omega)$、$\sigma_T(q,\omega)$ 的非局域结构。因此二维表面电导是三维双点介电核在严格零厚度极限下的自然接口，而不是人为引入的有效参数。

最后从同一边界方程求反射振幅。对 TE 波，切向电场连续给 $t_s=1+r_s$；磁场跳跃条件给
\[
Y_{s,1}(1-r_s)-Y_{s,2}t_s=\sigma_Tt_s.
\]
消去 $t_s$ 得
\[
r_s=\frac{Y_{s,1}-Y_{s,2}-\sigma_T}{Y_{s,1}+Y_{s,2}+\sigma_T}.
\]
对 TM 波，用标准 $p$ 极化单位矢量定义入射和反射幅；反射时基矢的切向分量随传播方向约定改变，边界代数给
\[
r_p=\frac{Y_{p,2}-Y_{p,1}+\sigma_L}{Y_{p,1}+Y_{p,2}+\sigma_L}.
\]
令 $\sigma_{L,T}\to0$，两式退化为普通 Fresnel 系数；令分母为零，则恢复正文表面本征模条件。

平面平移对称使散射 Green 张量对面内坐标差作二维 Fourier 对角化。源到界面再反射回观察点的法向传播相位是
\[
\ee^{\ii k_{z1}z'}\ee^{\ii k_{z1}z}=\ee^{\ii k_{z1}(z+z')}.
\]
自由一维 Helmholtz 预解式给 $\ii/(2k_{z1})$，再按 TE/TM 完备分解便得到正文\eqnrefs{eq:2d_sheet_reflected_green_weyl}。

对匀速平行电子，
\[
\int_{-\infty}^{\infty}\dd t\,\ee^{\ii\omega t-\ii q_xvt}
=2\pi\delta(\omega-q_xv),
\]
所以 $q_x=\omega/v$。若 $k_{z1}=\ii\kappa_1$ 且 $z=z'=b$，则
\[
\left|\ee^{\ii k_{z1}(z+z')}\right|=\ee^{-2\kappa_1b},
\]
得到正文\eqnrefs{eq:2d_eels_evanscent_impact_parameter_factor}。
\end{derivation}

\begin{derivation}[极性晶格的纵横光学频率关系]
从正文\eqnrefs{eq:polar-optical-oscillator-eom-dp76} 的受驱极性光学模式开始：
\[
M_r(\ddot u+\Gamma_{\rm ph}\dot u+\omega_{\rm TO}^2u)=Z^*eE,
\qquad
P_{\rm ion}=n_{\rm ion}Z^*eu.
\]
采用全文的 $\ee^{-\ii\omega t}$ 约定，时间导数映射为
\[
\dot u\mapsto-\ii\omega u(\omega),
\qquad
\ddot u\mapsto-\omega^2u(\omega).
\]
因此运动方程变成代数方程
\[
M_r\left(\omega_{\rm TO}^2-\omega^2-\ii\Gamma_{\rm ph}\omega\right)u(\omega)
=Z^*eE(\omega),
\]
解得
\[
u(\omega)
=\frac{Z^*e/M_r}
{\omega_{\rm TO}^2-\omega^2-\ii\Gamma_{\rm ph}\omega}E(\omega).
\]
离子极化于是为
\begin{align*}
P_{\rm ion}(\omega)
&=n_{\rm ion}Z^*eu(\omega)\\
&=\frac{n_{\rm ion}Z^{*2}e^2/M_r}
{\omega_{\rm TO}^2-\omega^2-\ii\Gamma_{\rm ph}\omega}E(\omega).
\end{align*}

高频电子极化已经吸收到背景介电常数 $\varepsilon_\infty$ 中，所以位移场为
\[
D(\omega)=\varepsilon_0\varepsilon_\infty E(\omega)+P_{\rm ion}(\omega)
=\varepsilon_0\varepsilon(\omega)E(\omega).
\]
除以 $\varepsilon_0E$：
\[
\varepsilon(\omega)
=\varepsilon_\infty
+\frac{n_{\rm ion}Z^{*2}e^2/(\varepsilon_0M_r)}
{\omega_{\rm TO}^2-\omega^2-\ii\Gamma_{\rm ph}\omega}.
\]
通分后
\begin{align*}
\varepsilon(\omega)
=\varepsilon_\infty
\frac{\omega_{\rm TO}^2-\omega^2-\ii\Gamma_{\rm ph}\omega
+n_{\rm ion}Z^{*2}e^2/(\varepsilon_0\varepsilon_\infty M_r)}
{\omega_{\rm TO}^2-\omega^2-\ii\Gamma_{\rm ph}\omega}.
\end{align*}
定义长波纵向光学频率
\[
\omega_{\rm LO}^2
\equiv\omega_{\rm TO}^2
+\frac{n_{\rm ion}Z^{*2}e^2}
{\varepsilon_0\varepsilon_\infty M_r},
\]
分子立即化为
$\omega_{\rm LO}^2-\omega^2-\ii\Gamma_{\rm ph}\omega$，所以
\[
\varepsilon(\omega)
=\varepsilon_\infty
\frac{\omega_{\rm LO}^2-\omega^2-\ii\Gamma_{\rm ph}\omega}
{\omega_{\rm TO}^2-\omega^2-\ii\Gamma_{\rm ph}\omega},
\]
即正文\eqnrefs{eq:polar_optical_dielectric_response}。

无阻尼时，$\omega=\omega_{\rm TO}$ 使分母为零，对应横向光学共振；$\omega=\omega_{\rm LO}$ 使 $\varepsilon(\omega)=0$，对应无外加自由电荷时纵向条件
$D=\varepsilon_0\varepsilon E=0$ 的非零电场解。静态极限再取 $\omega=0$、$\Gamma_{\rm ph}=0$：
\[
\varepsilon_s\equiv\varepsilon(0)
=\varepsilon_\infty
\frac{\omega_{\rm LO}^2}{\omega_{\rm TO}^2}.
\]
整理得到
\[
\frac{\varepsilon_s}{\varepsilon_\infty}
=\frac{\omega_{\rm LO}^2}{\omega_{\rm TO}^2},
\]
即正文\eqnrefs{eq:lattice_lst_relation}。
\end{derivation}

